\chapter{強化学習の基礎：MDP と価値関数}\label{ch:rl-basics}
\bookmap{6}
% ---- 格子世界の図を描くマクロ（数値は code/ch20/out/fig_*.tex をスクリプトが書き出す）
\providecommand{\rlbfont}{\footnotesize}
\providecommand{\rlbcell}[4]{\fill[#4] (#1,-#2) rectangle ++(1,-1);
  \node[font=\rlbfont,inner sep=0pt] at (#1+0.5,-#2-0.5) {$#3$};}
\providecommand{\rlbwall}[2]{\fill[black!50] (#1,-#2) rectangle ++(1,-1);}
\providecommand{\rlbgoal}[2]{\fill[bookgreen!75] (#1,-#2) rectangle ++(1,-1);
  \node[font=\rlbfont\bfseries,text=white,inner sep=0pt] at (#1+0.5,-#2-0.5) {G};}
\providecommand{\rlbpit}[2]{\fill[bookaccent!75] (#1,-#2) rectangle ++(1,-1);
  \node[font=\rlbfont\bfseries,text=white,inner sep=0pt] at (#1+0.5,-#2-0.5) {穴};}
\providecommand{\rlbarrow}[3]{\draw[->,very thick,bookblue]
  (#1+0.5,-#2-0.5) ++(#3:-0.3) -- ++(#3:0.6);}
\providecommand{\rlbframe}{\draw[black!45,step=1] (0,-3) grid (4,0);}
% 行動を選ぶ 9 状態の番号をマスの左上に小さく書く
\providecommand{\rlbnums}{\foreach \rlbs/\rlbc/\rlbr in
  {0/0/0,1/1/0,2/2/0,4/0/1,6/2/1,8/0/2,9/1/2,10/2/2,11/3/2}
  \node[font=\tiny,text=black!55,anchor=north west,inner sep=1pt]
    at (\rlbc,-\rlbr) {\rlbs};}
% 状態の番号に白い下地を付けた版（赤い枠や対角線が番号に重なる図で使う）
\providecommand{\rlbnumstag}{\foreach \rlbs/\rlbc/\rlbr in
  {0/0/0,1/1/0,2/2/0,4/0/1,6/2/1,8/0/2,9/1/2,10/2/2,11/3/2}
  \node[font=\tiny,text=black!60,fill=white,anchor=north west,inner sep=0.6pt]
    at (\rlbc,-\rlbr) {\rlbs};}
% その回に初めて値が入ったマスを囲む枠（価値反復のコマ送り）
\providecommand{\rlbnew}[2]{\draw[bookaccent,very thick]
  (#1+0.07,-#2-0.07) rectangle ++(0.86,-0.86);}
% 同じ軌跡を 3 回経験したときの MC と TD の表（\rlbft{行}{列}{値}{色}{変化}）
\providecommand{\rlbft}[5]{\fill[#4] (#2,-#1) rectangle ++(1,-1);
  \ifnum#5=1 \draw[bookaccent,very thick] (#2+0.06,-#1-0.06) rectangle ++(0.88,-0.88);\fi
  \node[font=\scriptsize,inner sep=0pt] at (#2+0.5,-#1-0.5) {#3};}
\providecommand{\rlbq}[6]{\begin{scope}[shift={(#1+0.5,-#2-0.5)}]
  \ifcase#3 \fill[#5] (0,0)--(-0.5,0.5)--(0.5,0.5)--cycle;
    \node[font=\rlbfont,inner sep=0pt] at (0,0.34) {\ifnum#6=1\bfseries\fi #4};
  \or \fill[#5] (0,0)--(0.5,0.5)--(0.5,-0.5)--cycle;
    \node[font=\rlbfont,inner sep=0pt] at (0.25,0) {\ifnum#6=1\bfseries\fi #4};
  \or \fill[#5] (0,0)--(0.5,-0.5)--(-0.5,-0.5)--cycle;
    \node[font=\rlbfont,inner sep=0pt] at (0,-0.34) {\ifnum#6=1\bfseries\fi #4};
  \or \fill[#5] (0,0)--(-0.5,-0.5)--(-0.5,0.5)--cycle;
    \node[font=\rlbfont,inner sep=0pt] at (-0.25,0) {\ifnum#6=1\bfseries\fi #4};
  \fi
  \draw[black!30,very thin] (-0.5,-0.5)--(0.5,0.5) (-0.5,0.5)--(0.5,-0.5);
  \end{scope}}
\begin{chapterlead}
\textbf{ここまで}: 第 V 部では、お手本をまねる学習（SFT）で、基盤モデルを計算アシスタントに作り変えました（\chapref{ch:posttraining}）。\\
\textbf{この章の問い}: 正解を教えてもらえず、結果の良し悪し（報酬）しか分からないとき、どう学習すればよいのでしょうか。
その数学的な土台は何でしょうか。\\
\textbf{この章を終えると}:
\begin{itemize}
\item 強化学習が教師あり学習と違う点を 3 つ挙げて説明できます。
\item 問題をマルコフ決定過程（MDP）として書き、方策・収益・価値関数を定義できます。
\item ベルマン方程式を 1 行ずつ導き、小さな格子世界で数値を確かめられます。
\item 価値反復、モンテカルロ法、TD 学習、Q 学習を NumPy で実装できます。
\item 表を使う方法が LLM に使えない理由と、次章の方策勾配法へ進む理由を説明できます。
\end{itemize}
\end{chapterlead}

\begin{readingguide}
\item[\rgpre] 第 VI 部（強化学習）の入口です。強化学習の知識は要りません。\chapref{ch:math}の確率・条件付き確率・期待値（\secref{sec:math-prob}）と、\chapref{ch:posttraining}の計算アシスタント（SFT）を思い出せれば読めます。微分と勾配降下法（\chapref{ch:ml}）を使うのは、流し読みできる小節だけです。
\item[\rgmust] \secref{sec:rlb-what}〜\secref{sec:rlb-value}：報酬、MDP、方策、収益、2 つの確率モデル（何を学習するのか）、価値関数、ベルマン期待方程式。次章は、ここで決めた言葉と記号をそのまま使います（\secref{sec:rlb-mdp}の小節「コードで書く」だけは流し読み可）。
\item[\rgmust] \secref{sec:rlb-dp}の小節「最適価値関数とベルマン最適方程式」と「価値反復」：平均を最大に変えると、最良の方策が求まります。Q 学習の土台です。
\item[\rgskim] \secref{sec:rlb-dp}の小節「方策改善と方策反復」：「評価と改善を交互に繰り返すと、最適方策に着く」という結論だけで先へ進めます。小節「床の滑りやすさを変えると、最適方策はどう変わるか」も、図を眺めるだけで足ります。
\item[\rglater] \secref{sec:rlb-dp}の小節「なぜ収束するのか」：縮小写像の紹介。あとの章では使いません。収束の理由が気になったときに戻れば足ります。
\item[\rgmust] \secref{sec:rlb-mc-td}〜\secref{sec:rlb-fa}：モンテカルロ法、TD 誤差、Q 学習、探索と活用、方策ベースの方法へのつながり。次章が使います（小節「実装して比べる」と「表は、いちばん単純なネットワーク」は流し読み可）。
\end{readingguide}

\chapref{ch:posttraining}で作った計算アシスタント（\file{sft.pt}）は、足し算と引き算の質問に 200 問中 87.0\% の割合で正しく答えました\footnote{\chapref{ch:posttraining}の本文が主に使ったのは、同じ設定で作り直した \file{ch16_sft.pt}（89.0\%）でした。ここから先は、共有のスクリプトで作った \file{sft.pt} を使います。200 問の正答率には数 \% のぶれがあり、2 つの違いはその範囲です（\chapref{ch:posttraining}のコラム「共有スクリプトの値との違い」）。}。
問題を 1,000 問に増やして測り直すと、正答率は 82.6\% です（\chapref{ch:rl-llm}では、こちらを出発点にします）。
このモデルは、私たちが用意した「質問と正しい答え」の組を、お手本としてまねて学習しました。

まねる学習には、限界が 2 つあります。
1 つ目は、\textbf{お手本の質が上限になる}ことです。
学習の目標が「お手本と同じ答えを出すこと」なので、お手本より上手になる理由がどこにもありません。
2 つ目は、\textbf{お手本をそもそも作れない}場合があることです。
道具を何十回も使い分けて進める仕事では、人間にも「正しい操作の一つ一つ」は書き下せません。

一方で、終わったあとに「うまくいったか」を判定することは、お手本を書くよりずっと簡単なことがよくあります。
計算問題なら、答えの数を照らし合わせるだけです。
この「結果の良し悪し」だけを手がかりに学ぶ方法が、\term{強化学習}{reinforcement learning, RL}です。

ただし、強化学習を使えば、どんな問題でも解けるわけではありません。
計算アシスタントの正答率がどこまで伸び、どこで伸び悩むのかは、\chapref{ch:rl-llm}と\chapref{ch:agentic-rl}で実際に試して数字で確かめます。
この章では、いったん LLM を離れます。
12 マスしかない小さな迷路を題材にして、強化学習の言葉と考え方を、手で計算できる大きさで身につけます。

この迷路には、ゴールと落とし穴があります。
でたらめに歩くと、ゴールに着けるのは 3 回に 1 回ほどです。
章の終わりには、迷路の仕組みを教えてもらえないエージェントが、スタートからゴールか穴に着くまでの試行を 2,000 回くり返すだけで、穴を避ける進み方を自分で見つけます。
教えるのは、「ゴールなら $+1$、穴なら $-1$」という結果の数だけです。
なぜそれだけで良い進み方が見つかるのかを、順に説明します。

章の前半（\secref{sec:rlb-dp}まで）は、迷路の仕組みがすべて分かっている場合を扱います。
後半は、仕組みを知らずに、歩いた経験だけから学ぶ場合です。
この章の方法をそのまま LLM に使うわけではありませんが、ここで出てくる言葉と考え方は、次章以降のすべての土台になります。

\section{強化学習とは何か}\label{sec:rlb-what}

\chapref{ch:ml}の学習方法の一覧で、強化学習は「正解は与えられず、答えの良し悪しの点数だけが返る学習」として一度だけ登場しました。
この節の問いは、「強化学習は教師あり学習と何が違い、なぜ LLM に必要なのか」です。

\subsection*{自転車の練習と、計算アシスタント}

自転車に乗る練習を思い出してください。
「ハンドルを右に 3 度切りなさい」と、一瞬ごとに正解を教えてくれる人はいません。
分かるのは、「転んだ」「前に進めた」という結果だけです。
それでも私たちは、試しては転ぶことを繰り返して、乗り方を身につけます。

計算アシスタントで言い直します。
SFT では、「12+35は？」という質問に「47です。」という正解の応答が付いていました。
強化学習では、正解の応答は与えられません。
モデルが自分で「57です。」と答えたあとで、「不正解」という判定だけが返ってきます。
何と答えるべきだったのかは、教えてもらえません。

\subsection*{教師あり学習との 3 つの違い}

教師あり学習との違いは、3 つあります（\tabref{tab:rlb-sl-vs-rl}）。

\textbf{(1) 正解ではなく報酬が返る。}
返ってくるのは、「どれくらい良かったか」を表す 1 つの数だけです。
この数を\term{報酬}{reward}と呼びます。

\textbf{(2) 報酬は遅れて届く。}
迷路では、ゴールに着いて初めて報酬がもらえます。
途中の何十歩のうち、どの一歩が良かったのかは分かりません。
功績をどの行動に割り振るかというこの問題を、\term{信用割当}{credit assignment}の問題と呼びます。

\textbf{(3) 自分の行動がデータを変える。}
教師あり学習では、データセットは学習の前に決まっていて、学習の途中で変わりません。
強化学習では、学習者の行動が、次に何を経験するかを決めます。
迷路の分かれ目でいつも右へ進んでいれば、左へ進んだ先に何があるのかは、いつまでも分かりません。
そのため、あえて試してみる\term{探索}{exploration}が必要になります（\secref{sec:rlb-q}）。

\begin{table}[tbp]
\centering
\caption{教師あり学習と強化学習の違い。}
\label{tab:rlb-sl-vs-rl}
\small
\begin{tabular}{lll}
\toprule
 & 教師あり学習（SFT など） & 強化学習 \\
\midrule
教えてもらえること & 正解そのもの & 結果の良し悪し（報酬） \\
いつ分かるか & 入力ごとにすぐ & 何歩もあとになることがある \\
データを決めるのは & 人間（データセット） & 学習者自身の行動 \\
到達できる上限 & お手本の質 & お手本を超えうる \\
計算アシスタントの例 & 「47です。」をまねる & 答えて、正誤だけを知る \\
\bottomrule
\end{tabular}
\end{table}

\subsection*{エージェントと環境}

強化学習では、問題に出てくるものを 2 つに分けます。
学習し、行動を選ぶ側が\term{エージェント}{agent}です。
エージェントの外側にあるすべてが\term{環境}{environment}です。
両者は、\figref{fig:rlb-agent-env}のように情報をやり取りします。

\begin{figure}[tbp]
\centering
\begin{tikzpicture}
  \node[box,minimum width=36mm,minimum height=12mm] (ag) at (0,0)
    {エージェント\\{\footnotesize 行動を選ぶ・学習する}};
  \node[greenbox,minimum width=36mm,minimum height=12mm] (env) at (7.4,0)
    {環境\\{\footnotesize 状態を変える・報酬を返す}};
  \draw[arr] ([yshift=2.5mm]ag.east) -- node[above,font=\small]{行動 $a_t$} ([yshift=2.5mm]env.west);
  \draw[arr] ([yshift=-2.5mm]env.west) -- node[below,font=\small]{報酬 $r_t$、次の状態 $s_{t+1}$}
    ([yshift=-2.5mm]ag.east);
\end{tikzpicture}
\caption{エージェントと環境のやり取り。時刻 $t$ に、エージェントは状態 $s_t$ を見て行動 $a_t$ を選ぶ。
環境は、報酬 $r_t$ と次の状態 $s_{t+1}$ を返す。返された $s_{t+1}$ が、次の時刻にエージェントが見る状態になる。これを繰り返す。
格子世界では、迷路を歩く人がエージェントで、迷路が環境。
計算アシスタントでは、LLM がエージェントで、環境は「質問を出し、LLM が書いたトークンを列の末尾に付け、書き終えたら採点する」仕組み（\secref{sec:rlb-fa}の\tabref{tab:rlb-maze-llm}）。}
\label{fig:rlb-agent-env}
\end{figure}

やり取りは、時刻 $t=0,1,2,\dots$ ごとに進みます。
エージェントは、いまの状況を表す\term{状態}{state} $s_t$ を受け取り、\term{行動}{action} $a_t$ を 1 つ選びます。
環境は、その結果として、報酬 $r_t$ と次の状態 $s_{t+1}$ を返します。
この 1 回のやり取りを、1 ステップと数えます。
迷路では 1 ステップが 1 歩に当たるので、この章では「1 歩」と書きます。
エージェントの目的は、受け取る報酬の合計をできるだけ大きくすることです（正確には「割引した合計の期待値」です。\secref{sec:rlb-return}で式にします）。

なお、\chapref{ch:agent}の「AI エージェント」は、道具を使って仕事を進める LLM です。
LLM をエージェント、道具や利用者を環境と見れば、これも\figref{fig:rlb-agent-env}の形に当てはまります。

\subsection*{なぜ LLM に強化学習が必要か}

LLM に強化学習を使う理由は、3 つあります。
1 つ目は、\textbf{お手本を超えたい}からです。
正誤さえ判定できれば、お手本にない解き方をモデルが自分で見つける余地があります。
2 つ目は、\textbf{正解が 1 つに決まらない}からです。
文章の良さに唯一の正解はなくても、「どちらの応答が良いか」という人間の好みなら集められます（以上、\chapref{ch:rl-llm}）。
3 つ目は、\textbf{長い手順の末に、やっと成否が分かる}からです。
道具を何度も使う仕事では、遅れて届く最後の結果しか採点できません（\chapref{ch:agentic-rl}）。

この節で分かったことをまとめます。
強化学習は、正解の代わりに報酬だけを受け取り、自分の行動で集めた経験から学ぶ方法です。
報酬が遅れて届くことと、行動がデータを変えることが、教師あり学習にはない難しさです。

\section{マルコフ決定過程：問題を数学の言葉で書く}\label{sec:rlb-mdp}

前節では、状態・行動・報酬という言葉を、日常の言葉として使いました。
この節の問いは、「環境の仕組みを、計算機で扱える形にどう書くか」です。

\subsection*{題材：格子世界}

この章では、\figref{fig:rlb-gridworld}の小さな迷路を最後まで使います。
強化学習の教科書で\jterm{格子世界}と呼ばれる、定番の題材です\cite{sutton2018rl}。
規則は次のとおりです。
\begin{itemize}
\item マスには、左上から順に 0 から 11 の番号を付けます。エージェントは、マス 8 から出発します。
\item マス 5 は壁で、入れません。
\item 行動は「上・右・下・左」の 4 つです。
\item 床が滑ります。確率 0.8 で選んだ向きに進みますが、確率 0.1 ずつで、その左右にそれます。左右は、選んだ向きから見た左手と右手です（「上」を選べば画面の左か右へ、「右」を選べば画面の上か下へそれます）。
\item 壁や外枠にぶつかったときは、その場にとどまります。
\item マス 3（ゴール）に入ると報酬 $+1$、マス 7（穴）に入ると報酬 $-1$ を受け取り、そこで終わります。それ以外の報酬は 0 です。
\end{itemize}

\begin{figure}[tbp]
\centering
\begin{tikzpicture}[x=1.02cm,y=1.02cm]
  \fill[bookgreen!75] (3,0) rectangle ++(1,-1);
  \fill[bookaccent!75] (3,-1) rectangle ++(1,-1);
  \fill[black!50] (1,-1) rectangle ++(1,-1);
  \fill[figblue] (0,-2) rectangle ++(1,-1);
  \rlbframe
  \foreach \s/\c/\r in {0/0/0,1/1/0,2/2/0,4/0/1,6/2/1,8/0/2,9/1/2,10/2/2,11/3/2}
    \node[font=\scriptsize,text=bookgray,anchor=north west,inner sep=2pt] at (\c,-\r) {\s};
  \foreach \s/\c/\r in {3/3/0,5/1/1,7/3/1}
    \node[font=\scriptsize,text=white,anchor=north west,inner sep=2pt] at (\c,-\r) {\s};
  \node[font=\footnotesize,text=white,align=center] at (3.55,-0.64) {ゴール\\$+1$};
  \node[font=\footnotesize,text=white,align=center] at (3.55,-1.64) {穴\\$-1$};
  \node[font=\small,text=white] at (1.5,-1.56) {壁};
  \node[font=\scriptsize] at (0.5,-2.6) {スタート};
  \begin{scope}[shift={(5.6,-1.9)}]
    \fill[bookblue] (0,0) circle (2pt);
    \draw[->,very thick,bookblue] (0,0) -- (0,1.0);
    \node[font=\footnotesize,align=center] at (0,1.35) {選んだ向き\\確率 0.8};
    \draw[->,thick,bookaccent] (0,0) -- (-0.7,0);
    \draw[->,thick,bookaccent] (0,0) -- (0.7,0);
    \node[font=\footnotesize,align=center] at (-1.15,0) {左手へ\\0.1};
    \node[font=\footnotesize,align=center] at (1.15,0) {右手へ\\0.1};
    \node[lbl] at (0,-0.5) {「上」を選んだとき};
  \end{scope}
  \begin{scope}[shift={(8.6,-1.5)}]
    \fill[bookblue] (0,0) circle (2pt);
    \draw[->,very thick,bookblue] (0,0) -- (1.0,0);
    \node[font=\footnotesize,align=center,anchor=west] at (1.05,0) {選んだ向き\\確率 0.8};
    \draw[->,thick,bookaccent] (0,0) -- (0,0.7);
    \draw[->,thick,bookaccent] (0,0) -- (0,-0.7);
    \node[font=\footnotesize] at (0,0.92) {左手へ 0.1};
    \node[font=\footnotesize] at (0,-0.92) {右手へ 0.1};
    \node[lbl] at (0.6,-1.4) {「右」を選んだとき};
  \end{scope}
\end{tikzpicture}
\caption{この章で使う格子世界。左上の小さな数字が状態の番号（3 がゴール、5 が壁、7 が穴、8 がスタート）。
床が滑るので、選んだ向きに進めるのは確率 0.8 で、確率 0.1 ずつで、選んだ向きから見た左手と右手にそれる（右の 2 つの図）。}
\label{fig:rlb-gridworld}
\end{figure}

床が滑るという規則は、わざと入れてあります。
同じ状態で同じ行動を選んでも、結果が毎回同じとは限りません。
現実の問題は、たいていそうなっています。

ゴールや穴のように、入ったら終わりになる状態を\term{終端状態}{terminal state}と呼びます。
スタートから終端状態に着くまでの 1 回分を、\term{エピソード}{episode}と呼びます。

\subsection*{MDP の 5 つの部品}

格子世界の規則を、記号で書き直します。
まず、状態 8 で「右」を選んだとします。
確率 0.8 で、右隣の状態 9 に進みます。
確率 0.1 で、進む向きの左手（上）にそれて、状態 4 に入ります。
確率 0.1 で、右手（下）にそれますが、そこは外枠なので、状態 8 にとどまります。
これを、条件付き確率（\chapref{ch:math}）の記号で
\[
P(9\mid 8,\text{右})=0.8,\qquad P(4\mid 8,\text{右})=0.1,\qquad P(8\mid 8,\text{右})=0.1
\]
と書きます。
縦棒の右側が条件（いまの状態と選んだ行動）で、左側が次の状態です。
この $P$ が、すぐあとの定義に出てくる遷移確率です。

この関係を絵にしたものが、\figref{fig:rlb-mdp-graph}の\jterm{状態遷移図}です。
丸が状態、黒い点が「状態と行動の組」です。
状態 8 の丸から出る青い矢印が、エージェントが「右」を選ぶことを表します。
黒い点から出る橙の 3 本の矢印が、そのあと環境が決める行き先で、矢印に添えた数（0.8、0.1、0.1）が確率です。
3 本の確率を足すと 1 になります。
状態 6 で「上」を選んだ場合も描きました。
こちらは、確率 0.1 で穴に入り、報酬 $-1$ を受け取ります。

\begin{figure}[tbp]
\centering
\begin{tikzpicture}[x=1.0cm,y=1.0cm,
  st/.style={circle,draw=black!55,fill=white,minimum size=7mm,inner sep=0pt,font=\small},
  ac/.style={circle,fill=black!80,minimum size=2.4mm,inner sep=0pt},
  pol/.style={->,very thick,bookblue},
  env/.style={->,thick,bookorange},
  pr/.style={font=\footnotesize,fill=white,inner sep=1pt}]
  \foreach \n/\x/\y in {0/0/0,1/2.7/0,2/5.4/0,4/0/-2.1,6/5.4/-2.1,8/0/-4.2,9/2.7/-4.2,10/5.4/-4.2,11/8.1/-4.2}
    \node[st] (n\n) at (\x,\y) {\n};
  \node[st,fill=bookgreen!75,text=white,font=\small\bfseries] (n3) at (8.1,0) {3};
  \node[font=\scriptsize,anchor=west] at (8.5,0) {ゴール};
  \node[st,fill=bookaccent!75,text=white,font=\small\bfseries] (n7) at (8.1,-2.1) {7};
  \node[font=\scriptsize,anchor=west] at (8.5,-2.1) {穴};
  \fill[black!45] (2.4,-2.4) rectangle (3.0,-1.8);
  \node[font=\scriptsize,text=white] at (2.7,-2.1) {壁};
  \node[font=\scriptsize,anchor=north] at (0,-4.6) {スタート};
  % 状態 8 で「右」
  \node[ac] (a8) at (1.25,-3.15) {};
  \draw[pol] (n8) -- node[pos=0.45,above left=-1pt,font=\footnotesize,text=bookblue]{右} (a8);
  \draw[env] (a8) -- node[pos=0.5,above right=-1pt,font=\footnotesize,text=bookorange]{0.8} (n9);
  \draw[env] (a8) -- node[pos=0.5,above right=-1pt,font=\footnotesize,text=bookorange]{0.1} (n4);
  \draw[env] (a8) to[bend left=60,looseness=1.3] node[pos=0.5,below right=-1pt,font=\footnotesize,text=bookorange]{0.1} (n8);
  % 状態 6 で「上」
  \node[ac] (a6) at (6.9,-0.85) {};
  \draw[pol] (n6) -- node[pos=0.5,above left=-1pt,font=\footnotesize,text=bookblue]{上} (a6);
  \draw[env] (a6) -- node[pos=0.5,above right=-1pt,font=\footnotesize,text=bookorange]{0.8} (n2);
  \draw[env] (a6) -- node[pos=0.45,right=2pt,font=\footnotesize,text=bookorange]{0.1、$r=-1$} (n7);
  \draw[env] (a6) to[bend left=60,looseness=1.3] node[pos=0.5,below right=-1pt,font=\footnotesize,text=bookorange]{0.1} (n6);
  % 凡例
  \draw[pol] (2.0,-5.3) -- (2.6,-5.3);
  \node[font=\footnotesize,anchor=west] at (2.7,-5.3) {エージェントが選ぶ（方策 $\pi$）};
  \draw[env] (7.0,-5.3) -- (7.6,-5.3);
  \node[font=\footnotesize,anchor=west] at (7.7,-5.3) {環境が決める（遷移確率 $P$）};
\end{tikzpicture}
\caption{格子世界の状態遷移図（一部）。丸が状態、黒い点が「状態と行動の組」。
青い矢印はエージェントが行動を選ぶところ、橙の矢印は環境が次の状態を決めるところで、矢印に添えた数は遷移確率 $P(s'\mid s,a)$。
報酬は 0 でないものだけを書いた。
描いたのは、状態 8 で「右」と、状態 6 で「上」の 2 組だけ（9 状態 $\times$ 4 行動の 36 組を全部描くと、黒い点が 36 個になる）。
自分に戻る矢印は、壁や外枠にぶつかってとどまる場合。}
\label{fig:rlb-mdp-graph}
\end{figure}

図の青と橙は、役割の違いを表しています。
エージェントが決められるのは、青い矢印（どの向きを選ぶか）だけです。
選んだあとにどの橙の矢印が起きるかは、環境が確率で決めます。
エージェントは、行き先のマスを直接は選べません。
この書き方で環境全体を表したものが、マルコフ決定過程です。

\begin{definition}[マルコフ決定過程]\label{def:rlb-mdp}
状態集合 $\calS$、行動集合 $\calA$、遷移確率 $P$、報酬関数 $r$、割引率 $\gamma$ の 5 つの組 $(\calS,\calA,P,r,\gamma)$ を、\term{マルコフ決定過程}{Markov decision process, MDP}と呼びます。
\jterm{遷移確率} $P(s'\mid s,a)$ は、状態 $s$ で行動 $a$ を選んだとき、次の状態が $s'$ になる確率です。
\jterm{報酬関数} $r(s,a,s')$ は、そのとき受け取る報酬です。
割引率 $\gamma$ は 0 以上 1 以下の数で、将来の報酬をどれだけ割り引くかを決めます（\secref{sec:rlb-return}）。
\end{definition}

$s'$ は「次の状態」を表す記号で、「エス・プライム」と読みます。
報酬が次の状態によらない問題では、報酬関数を $r(s,a)$ と短く書きます。
次章では、簡単のために、この短い形を使います。
5 つの組に加えて、最初の状態 $s_0$ を決める確率分布 $p(s_0)$ も、問題の一部です。
格子世界では、$s_0$ はいつも 8 です。
MDP に入っているのは、環境の側の決まりだけです。
エージェントの振る舞い（方策）は MDP の外にあり、次節で別に決めます。
\tabref{tab:rlb-mdp}に、格子世界との対応をまとめます。

\begin{table}[tbp]
\centering
\caption{MDP の 5 つの部品と、格子世界での中身。}
\label{tab:rlb-mdp}
\small
\begin{tabular}{lll}
\toprule
記号 & 名前 & 格子世界では \\
\midrule
$\calS$ & 状態集合 & マスの番号（壁を除く 11 個。終端 2 個を除く 9 個で行動を選ぶ） \\
$\calA$ & 行動集合 & 上・右・下・左の 4 個 \\
$P(s'\mid s,a)$ & 遷移確率 & 確率 0.8 で選んだ向き、0.1 ずつでその左手と右手 \\
$r(s,a,s')$ & 報酬関数 & $s'$ がゴールなら $+1$、穴なら $-1$、ほかは 0 \\
$\gamma$ & 割引率 & 0.9（この章を通して固定） \\
\bottomrule
\end{tabular}
\end{table}

\subsection*{マルコフ性}

MDP の「マルコフ」は、次の性質を指しています。

\begin{definition}[マルコフ性]\label{def:rlb-markov}
次の状態と報酬の確率が、\textbf{いまの状態と行動だけ}で決まり、それより前の経緯によらないことを、\term{マルコフ性}{Markov property}と呼びます。
式で書くと、次のとおりです。
\[
P(s_{t+1}\mid s_t,a_t,\,s_{t-1},a_{t-1},\dots,s_0,a_0)=P(s_{t+1}\mid s_t,a_t)
\]
報酬は $r(s_t,a_t,s_{t+1})$ で決まるので、式には次の状態だけを書いた。
\end{definition}

格子世界は、この性質を満たします。
次にどのマスへ進むかは、いまいるマスと選んだ向きだけで決まります。
どの経路でそのマスに来たかは、関係ありません。

\begin{intuition}{状態は「これまでの要約」}
将棋では、いまの盤面と持ち駒さえ分かれば、次の一手を考えられます。
それまでの指し手の順番を覚えている必要はありません。
マルコフ性とは、「先のことを考えるのに必要な情報が、状態にすべて入っている」ということです。
入っていない場合は、状態に含める情報を増やします。
言語モデルでは、「ここまでのトークン列の全体」を状態にします（\chapref{ch:rl-llm}）。
\end{intuition}

マルコフ性があると、「この状態から先がどれくらい良いか」を、状態ごとに 1 つの数で表せます。
これが、\secref{sec:rlb-value}の価値関数です。

\subsection*{コードで書く}
\skimread{押さえる結論は 3 つです。表 \code{P[s][a]} には、「（確率, 次の状態, 報酬）」の組のリストが入ります。\code{STATES} は、行動を選ぶ 9 個の状態の番号です。エージェントが環境に対して呼べる関数は、\code{reset} と \code{step} だけです。}

格子世界をコードにします。
\coderef{code:rlb-grid-def}は、定数の定義です。

\snippet{格子世界の定数。状態と行動を整数で表す}{code:rlb-grid-def}{ch20/gridworld.py}{grid-def}

状態は、「$4\times$行$+$列」という 1 つの整数で表します（行と列は 0 から数えます）。
マスは 12 個、壁を除いた状態は 11 個で、そのうちエージェントが行動を選ぶのは、終端状態を除く 9 個（\code{STATES}）です。
この先のコードでは、マスの番号をそのまま配列の添字に使うので、長さ 12 の配列を用意し、壁と終端状態の欄は 0 のまま使いません。
行動は時計回りに 0 から 3 の番号を付けました。
こうしておくと、「左にそれる」（進む向きから見た左手へそれる）は番号を 1 減らすこと、「右にそれる」は 1 増やすことに対応します。

\coderef{code:rlb-grid-model}は、遷移確率と報酬の表を作ります。

\snippet{遷移確率と報酬の表 \protect\code{P} を作る}{code:rlb-grid-model}{ch20/gridworld.py}{grid-build}

\code{build_model} は、すべての状態と行動について、起こりうる結果を並べます。
8〜10 行目が、確率 0.8 の「選んだ向き」と、確率 0.1 ずつの「左右」です。
行動の番号に 1 を足し引きしたあと、4 で割った余りを取って、番号を一周させています（Python では、$-1$ を 4 で割った余りは 3 です）。
11 行目の \code{move(s, d)} は、状態 \code{s} から向き \code{d} に 1 マス進んだ先を返す関数です。
壁か外枠にぶつかるときは、\code{s} をそのまま返します。
12 行目では、同じ状態に行き着く場合の確率を足し合わせています。
たとえば、状態 8 で「下」を選ぶと、下へ進む場合（0.8）も、進む向きの右手、つまり画面の左へそれる場合（0.1）も、外枠にぶつかります。
この 2 つが足されて、「確率 0.9 で状態 8 にとどまる」になります。
残りの 0.1（左手、つまり画面の右へそれる場合）だけが、状態 9 に進みます。
13 行目の \code{reward(s2)} は、入った先がゴールなら $+1$、穴なら $-1$、ほかは 0 を返します。
結果は、\code{P[s][a]} に「（確率, 次の状態, 報酬）」の組のリストとして入ります。
中身を 2 つ見てみます。

\outputfile{ch20/out/grid_model.txt}

状態 8 で「右」の結果は、先ほど手で書いた $P(\cdot\mid 8,\text{右})$ と同じです。
状態 6 で「上」を選ぶと、確率 0.1 で右手にそれて穴に落ちます（\figref{fig:rlb-mdp-graph}の右上）。

現実の問題では、エージェントはこの表を見ることができません。
使えるのは、行動を渡すと次の状態と報酬を返す関数だけです。
\coderef{code:rlb-grid-env}が、その関数を持つクラスです。

\snippet{環境のクラス。エージェントは \protect\code{reset} と \protect\code{step} だけを使う}{code:rlb-grid-env}{ch20/gridworld.py}{grid-env}

\code{reset} は、エージェントをスタートに戻します。
\code{step} は、行動 \code{a} を受け取り、表 \code{P} の確率に従って結果を 1 つ選びます。
13 行目の \code{zip(*...)} は、組のリストを、確率・次の状態・報酬の 3 本に分け直します。
14 行目の \code{rng.choice} は、確率 \code{probs} に従って番号を 1 つ選びます。
返すのは、次の状態、報酬、終わったかどうかの 3 つです。

この節で分かったことをまとめます。
環境は、状態・行動・遷移確率・報酬関数・割引率の 5 つで書けます。
マルコフ性により、先のことは、いまの状態だけから考えられます。

\section{方策・軌跡・収益：良さを測る基準}\label{sec:rlb-return}

環境を書けたので、次はエージェントの側です。
この節の問いは、「エージェントの振る舞いをどう書き、その良し悪しを何で測るか」です。

\subsection*{方策}

エージェントの振る舞いとは、「どの状態で、どの行動を選ぶか」の決まりです。
この決まりを\term{方策}{policy}と呼び、$\pi$（パイ）で表します。
$\pi(a\mid s)$ は、状態 $s$ で行動 $a$ を選ぶ確率です。
たとえば、どの状態でも 4 つの行動を等しい確率で選ぶ\jterm{ランダム方策}は、すべての $s$ と $a$ について $\pi(a\mid s)=1/4$ です。
確率で行動を選ぶ方策を\jterm{確率的方策}、各状態で行動を 1 つに決めている方策を\jterm{決定的方策}と呼びます。
決定的方策では、状態 $s$ で選ぶ行動を $\pi(s)$ と書きます。

方策の形に、見覚えはないでしょうか。
言語モデルは、ここまでのトークン列を受け取って、次のトークンの確率を出す関数でした（\chapref{ch:lm}）。
「状態を受け取って、行動の確率を出す」という方策と、同じ形をしています（章の最後で、もう一度取り上げます）。

\subsection*{軌跡}

方策 $\pi$ に従って環境とやり取りすると、状態・行動・報酬が順に並んだ記録ができます。
\begin{equation}
\tau=(s_0,a_0,r_0,\ s_1,a_1,r_1,\ \dots,\ s_{T-1},a_{T-1},r_{T-1})
\label{eq:rlb-trajectory}
\end{equation}
これを\term{軌跡}{trajectory}と呼び、$\tau$（タウ）で表します。
$T$ はエピソードの長さ（行動した回数）です。
各時刻で起きることを、順に書きます。
「$\sim$」は、「右側の確率分布から 1 つ選ぶ」という意味でした（\chapref{ch:math}）。
\begin{align*}
s_0&\sim p(s_0) && \text{最初の状態が決まる}\\
a_t&\sim \pi(\cdot\mid s_t) && \text{エージェントが方策に従って行動を選ぶ}\\
s_{t+1}&\sim P(\cdot\mid s_t,a_t) && \text{環境が次の状態を決める}\\
r_t&=r(s_t,a_t,s_{t+1}) && \text{報酬が決まる}
\end{align*}
本書では、行動 $a_t$ の直後に受け取る報酬を $r_t$ と書きます\footnote{同じ報酬を $r_{t+1}$ と書く教科書もあります\cite{sutton2018rl}。添字のずれに注意してください。}。

4 行のうち、2 行目と 3 行目に確率が出てきます。
2 行目の $\pi$ はエージェントの選び方で、3 行目の $P$ は環境の規則です。
確率の連鎖律（\chapref{ch:math}）で各時刻の確率を順に掛けると、1 本の軌跡 $\tau$ が起きる確率になります。
各時刻の確率がいまの状態と行動だけで決まる（マルコフ性）ので、掛ける因子は短く書けます。
\begin{equation}
p_\pi(\tau)=p(s_0)\prod_{t=0}^{T-1}\underbrace{\pi(a_t\mid s_t)}_{\text{方策}}\;\underbrace{P(s_{t+1}\mid s_t,a_t)}_{\text{環境の遷移確率}}
\label{eq:rlb-traj-prob}
\end{equation}
$p_\pi$ の添字 $\pi$ は、「方策 $\pi$ で動いたときの」軌跡の確率であることを示します。
報酬は $s_t,a_t,s_{t+1}$ から決まるので、確率の式には現れません。

\coderef{code:rlb-rollout}は、1 エピソード分の軌跡を作る関数です。

\snippet{方策に従って 1 エピソード動かし、収益を計算する}{code:rlb-rollout}{ch20/gridworld.py}{grid-rollout}

引数の \code{policy} は、状態ごとに 4 つの行動の確率を並べたものです。
\code{run_episode} の 5 行目で、この確率に従って行動を選びます。
6 行目で、環境が次の状態と報酬を返します。
7 行目で、（状態, 行動, 報酬）の組を記録します。
\code{returns} は、すぐあとで説明する「収益」を計算します。
ランダム方策で 1 エピソード動かした結果が、\figref{fig:rlb-trajectory}です。

\begin{figure}[tbp]
\centering
\begin{tikzpicture}[x=1.05cm,y=1.05cm]
  \fill[bookgreen!75] (3,0) rectangle ++(1,-1);
  \fill[bookaccent!75] (3,-1) rectangle ++(1,-1);
  \fill[black!50] (1,-1) rectangle ++(1,-1);
  \rlbframe
  \foreach \s/\c/\r in {0/0/0,1/1/0,2/2/0,4/0/1,6/2/1,8/0/2,9/1/2,10/2/2,11/3/2}
    \node[font=\scriptsize,text=bookgray,anchor=north west,inner sep=2pt] at (\c,-\r) {\s};
  \foreach \s/\c/\r in {3/3/0,5/1/1,7/3/1}
    \node[font=\scriptsize,text=white,anchor=north west,inner sep=2pt] at (\c,-\r) {\s};
  \node[font=\footnotesize\bfseries,text=white] at (3.55,-0.8) {ゴール};
  \node[font=\footnotesize\bfseries,text=white] at (3.5,-1.5) {穴};
  \node[font=\footnotesize,text=white] at (1.5,-1.5) {壁};
  \tikzset{rlbt/.style={midway,circle,draw,fill=white,inner sep=0.7pt,font=\scriptsize}}
  \draw[->,very thick,bookblue] (0.5,-2.5) -- node[rlbt]{0} (1.4,-2.5);
  \draw[->,very thick,bookblue] (1.5,-2.5) -- node[rlbt]{1} (2.4,-2.5);
  \draw[->,very thick,bookaccent] (2.5,-2.5) -- node[rlbt,pos=0.7]{2} (2.5,-1.6);
  \draw[->,very thick,bookblue] (2.5,-1.5) -- node[rlbt]{3} (2.5,-0.6);
  \draw[->,very thick,bookblue] (2.5,-0.5) -- node[rlbt]{4} (3.4,-0.5);
  \node[lbl,anchor=north west,align=left] at (0,-3.05) {矢印の数字は時刻 $t$。赤い矢印（$t=2$）では、\\「右」を選んだが、滑って上に進んだ};
  \node[anchor=north west,font=\footnotesize] at (4.3,0.1) {%
    \begin{tabular}{@{}ccccc@{}}
    \toprule
    $t$ & $s_t$ & $a_t$ & $r_t$ & $G_t$ \\
    \midrule
    0 & 8 & 右 & 0 & 0.6561 \\
    1 & 9 & 右 & 0 & 0.7290 \\
    2 & 10 & 右 & 0 & 0.8100 \\
    3 & 6 & 上 & 0 & 0.9000 \\
    4 & 2 & 右 & $+1$ & 1.0000 \\
    \bottomrule
    \end{tabular}};
\end{tikzpicture}
\caption{ランダム方策で動かした軌跡の例（\file{ch20/gridworld.py} の実行結果）。長さは $T=5$。
説明しやすいように、短く終わった回を選んだ（平均の長さは約 32 歩）。
報酬は最後の 1 歩でだけ届く。$G_t$ は、このあと定義する収益（$\gamma=0.9$）。}
\label{fig:rlb-trajectory}
\end{figure}

報酬は最後の 1 歩でだけ届きます。
\secref{sec:rlb-what}の「遅れて届く報酬」の簡単な例です。

\subsection*{収益と割引}

エージェントの目的は、報酬の合計を大きくすることでした。
ただし、単純な合計ではなく、遠い将来の報酬ほど小さく数えた合計を使います。

\begin{definition}[収益]\label{def:rlb-return}
時刻 $t$ から先に受け取る報酬を、1 歩先になるごとに $\gamma$ 倍に割り引いて足したものを、\term{収益}{return}と呼びます。
\begin{equation}
G_t\defeq r_t+\gamma r_{t+1}+\gamma^2 r_{t+2}+\cdots+\gamma^{T-1-t}r_{T-1}
=\sum_{k=0}^{T-1-t}\gamma^k r_{t+k}
\label{eq:rlb-return}
\end{equation}
$\gamma$ を\term{割引率}{discount factor}と呼びます。
軌跡全体の収益 $G_0$ を、$R(\tau)$ とも書きます。
\end{definition}

\figref{fig:rlb-trajectory}の軌跡で計算してみます。
報酬は $(r_0,\dots,r_4)=(0,0,0,0,1)$ で、$\gamma=0.9$ です。
\[
G_0=0+0.9\cdot0+0.9^2\cdot0+0.9^3\cdot0+0.9^4\cdot1=0.6561
\]
表のいちばん上の行と一致します。
同じゴールでも、たどり着くのが 4 歩先なので、$+1$ が $0.9^4\approx0.66$ に割り引かれています。

なぜ割り引くのでしょうか。理由は 3 つあります。
\begin{itemize}
\item \textbf{早く着くほうを良いとするため。}
$\gamma=1$ だと、5 歩で着いても 50 歩で着いても、収益は同じ $+1$ です。
$\gamma=0.9$ なら、5 歩なら $0.9^4\approx0.66$、50 歩なら $0.9^{49}\approx0.006$ で、近道のほうがはっきり良くなります。
\item \textbf{合計が無限に大きくならないようにするため。}
終わりのない問題でも、$\gamma<1$ なら、報酬が毎回 1 以下のとき、収益は $1+\gamma+\gamma^2+\cdots=1/(1-\gamma)$ 以下に収まります（この等式は、すぐあとで確かめます）。
\item \textbf{遠い将来は不確かだから。}
先のことほど予想しにくいので、控えめに数えます。
\end{itemize}

2 つ目の理由で使った等式 $1+\gamma+\gamma^2+\cdots=1/(1-\gamma)$ を確かめます。
$\gamma<1$ なら、項は先へ行くほど小さくなり、和は決まった値に落ち着きます。
その値を $S$ と置いて、全体に $\gamma$ を掛けると、$\gamma S=\gamma+\gamma^2+\gamma^3+\cdots$ です。
右辺は、$S$ から最初の 1 を除いたものなので、$\gamma S=S-1$ です。
これを $S$ について解くと、$S=1/(1-\gamma)$ になります。
$\gamma=0.9$ なら $S=1/0.1=10$、$\gamma=0.99$ なら $S=100$ です。

\figref{fig:rlb-discount}は、$k$ 歩先の報酬にかかる重み $\gamma^k$ です。
重みの合計は、いま求めた $1/(1-\gamma)$ なので、おおまかには「$1/(1-\gamma)$ 歩先までを数えている」と考えられます。
$\gamma=0.9$ なら 10 歩先、$\gamma=0.99$ なら 100 歩先までが目安です。

\begin{figure}[tbp]
\centering
\begin{tikzpicture}
\begin{axis}[height=0.32\linewidth,xlabel={何歩先の報酬か $k$},ylabel={重み $\gamma^k$},
  xmin=0,xmax=40,ymin=0,ymax=1.05,legend pos=north east]
\addplot[bookaccent,domain=0:40,samples=41] {0.5^x};
\addplot[bookblue,domain=0:40,samples=41] {0.9^x};
\addplot[bookgreen,domain=0:40,samples=41] {0.99^x};
\legend{$\gamma=0.5$,$\gamma=0.9$,$\gamma=0.99$}
\end{axis}
\end{tikzpicture}
\caption{割引の重み $\gamma^k$。$\gamma$ が小さいほど、目先の報酬だけを見る。
$\gamma$ が 1 に近いほど、遠い将来の報酬まで数える。}
\label{fig:rlb-discount}
\end{figure}

\subsection*{収益の再帰式}

収益には、この章で何度も使う大事な性質があります。
式\eqref{eq:rlb-return}の 2 項目以降を、$\gamma$ でくくります。
\begin{align}
G_t&=r_t+\gamma r_{t+1}+\gamma^2 r_{t+2}+\cdots \notag\\
&=r_t+\gamma\,(r_{t+1}+\gamma r_{t+2}+\cdots) && \text{2 項目以降を $\gamma$ でくくる}\notag\\
&=r_t+\gamma\,G_{t+1} && \text{かっこの中は、時刻 $t+1$ からの収益}
\label{eq:rlb-return-recursive}
\end{align}
「いまからの収益」は、「すぐにもらう報酬」と「次の時刻からの収益の $\gamma$ 倍」の和です。
最後の行動のあとには報酬がないので、$G_T=0$ と決めておきます。
\coderef{code:rlb-rollout}の \code{returns} は、この式を後ろから前へ順に使っています。
\figref{fig:rlb-trajectory}の表でも、$G_4=1$、$G_3=0+0.9\times1=0.9$、$G_2=0.81$ と、下から順に計算できます。

\begin{quiz}
$\gamma=0.5$ で、報酬が $(r_0,r_1,r_2)=(1,0,8)$ の 3 歩で終わる軌跡があります。
$G_2$、$G_1$、$G_0$ を、式\eqref{eq:rlb-return-recursive}で後ろから順に求めてください。
\end{quiz}

答えは、$G_2=8$、$G_1=0+0.5\times8=4$、$G_0=1+0.5\times4=3$ です。

同じ方策でも、収益は毎回違います。
床が滑り、方策も確率で行動を選ぶので、軌跡が毎回変わるからです。
ランダム方策で 1 万エピソードを動かしてみました（\file{ch20/gridworld.py} の後半）。
ゴールに着けたのは 35.3\% で、残りは穴に落ちました。
収益 $G_0$ はエピソードごとに正にも負にもなり、平均は $-0.066$ でした。
方策の良さは、1 回の収益ではなく、\textbf{収益の期待値}（\chapref{ch:math}）で測るのが自然です。
強化学習の目的は、収益の期待値がいちばん大きくなる方策を見つけることです。
式で書くと、スタートからの収益 $G_0$ の期待値を、方策の良さとします。
\begin{equation}
J(\pi)\defeq\E_{\tau\sim p_\pi}[\,G_0\,]
\label{eq:rlb-objective}
\end{equation}
$J$ を\term{期待収益}{expected return}と呼びます。
期待値は、式\eqref{eq:rlb-traj-prob}の確率 $p_\pi(\tau)$ を重みにした、軌跡についての平均です。
ランダム方策では、1 万エピソードの平均から $J\approx-0.066$ です。
$J$ がいちばん大きい方策を探すことが、強化学習の問題です。

\subsection*{確率モデルは 2 つ：何を学習するのか}

ここまでで、この章に出てくる確率モデルが出そろいました。
2 つあります（\tabref{tab:rlb-two-models}）。
\begin{itemize}
\item \textbf{環境の遷移確率 $P(s'\mid s,a)$}。迷路の規則（床の滑り方）です。問題として与えられたもので、学習では変わりません。エージェントには変えられません。
\item \textbf{方策 $\pi(a\mid s)$}。各マスで、どの向きをどの確率で選ぶかです。これを学習で良くします。
\end{itemize}
軌跡の確率（式\eqref{eq:rlb-traj-prob}）には、この 2 つの因子が交互に並んでいます。
学習で動かせるのは、$\pi$ の因子だけです。

\begin{table}[tbp]
\centering
\caption{この章に出てくる 2 つの確率モデル。学習で良くするのは方策だけで、遷移確率は与えられたまま変わらない。}
\label{tab:rlb-two-models}
\small
\begin{tabular}{lll}
\toprule
 & 環境の遷移確率 $P(s'\mid s,a)$ & 方策 $\pi(a\mid s)$ \\
\midrule
何の確率か & 行動のあと、次にどの状態になるか & 状態 $s$ で、どの行動を選ぶか \\
誰のものか & 環境（MDP の一部） & エージェント \\
学習で変わるか & 変わらない（与えられたもの） & 変わる（これを良くする） \\
格子世界では & 床の滑り（0.8、0.1、0.1） & 各マスでの向きの選び方 \\
LLM では & トークンを末尾に付ける（決定的） & 言語モデルそのもの（\chapref{ch:rl-llm}） \\
\bottomrule
\end{tabular}
\end{table}


では、「スタートから、どの確率でどのマスへ進めばよいか、というマスの列の確率を学習する」という理解は正しいでしょうか。
これは、半分正しい理解です。
正しいのは、学習の結果として、マスの列（軌跡）の確率が変わることです。
方策を決めると、式\eqref{eq:rlb-traj-prob}によって、どの軌跡がどの確率で起きるかが決まります。
たとえば\figref{fig:rlb-trajectory}の軌跡は、ランダム方策では $(1/4)^5\times0.8^4\times0.1\approx4\times10^{-5}$ の確率で起きます（5 回の行動がそれぞれ $1/4$、遷移は 0.8 が 4 回と、滑った 0.1 が 1 回）。
最適方策（\secref{sec:rlb-dp}）は状態 8 で「右」を選ばないので、同じ軌跡の確率は 0 です。

正しくないのは、「どのマスに行くか」を直接学ぶと考える点です。
学習で決めるのは、各マスでの向きの選び方 $\pi(a\mid s)$ です。
行き先のマスは、そのあと環境の $P$ が確率で決めます。
また、マスの列そのものの確率を、表として覚えるわけでもありません。
列の確率は、学んだ $\pi$ と、与えられた $P$ の積として、あとから決まります。

方策を 1 つに決めると、行動の選び方の確率と遷移確率を合わせて、「マスからマスへ移る確率」
\[
P_\pi(s'\mid s)=\sum_a\pi(a\mid s)\,P(s'\mid s,a)
\]
だけが残ります。
行動を挟まずに、状態だけが確率で移っていく過程になるわけです。
これを\term{マルコフ連鎖}{Markov chain}と呼びます。
\figref{fig:rlb-chain}は、ランダム方策と最適方策のそれぞれで、この確率を矢印の太さで描いたものです（\file{ch20/fig_dp_extra.py} で計算）。
左のランダム方策では、どの矢印も細く（確率 0.25）、穴へ向かう矢印も残っています。
右の最適方策では、ゴールへ向かう太い矢印（確率 0.8）の列ができ、穴へ入る矢印は状態 6 と 11 からの 0.1 だけです。
方策を変えると、マスからマスへの確率が変わり、その結果として、マスの列の確率が変わります。
これが、「半分正しい」の意味です。

\begin{figure}[tbp]
\centering
\begin{tikzpicture}[x=1.15cm,y=1.05cm,
  cs/.style={circle,draw=black!55,fill=white,minimum size=5.5mm,inner sep=0pt,font=\scriptsize}]
  \foreach \k/\fn/\ttl in {0/random/{ランダム方策（どの行動も $1/4$）},1/star/{最適方策 $\pi^\ast$（\secref{sec:rlb-dp}）}} {
  \begin{scope}[shift={(5.6*\k,0)}]
    \node[font=\small] at (1.5,0.75) {\ttl};
    \foreach \n/\c/\r in {0/0/0,1/1/0,2/2/0,4/0/1,6/2/1,8/0/2,9/1/2,10/2/2,11/3/2}
      \node[cs] (c\n) at (\c,-\r) {\n};
    \node[cs,fill=bookgreen!75,text=white] (c3) at (3,0) {G};
    \node[cs,fill=bookaccent!75,text=white] (c7) at (3,-1) {穴};
    \fill[black!45] (0.75,-1.25) rectangle (1.25,-0.75);
    \input{code/ch20/out/fig_chain_\fn.tex}
  \end{scope}}
\end{tikzpicture}
\caption{方策を決めると、マスからマスへ移る確率 $P_\pi(s'\mid s)=\sum_a\pi(a\mid s)P(s'\mid s,a)$ が決まる。矢印の太さが確率（最も太いものが 0.8、最も細いものが 0.1）。
その場にとどまる確率は描いていない。G はゴール（状態 3）、穴は状態 7、灰色の四角は壁。
左のランダム方策では、どの矢印も確率 0.25。右の最適方策では、ゴールへ向かう確率 0.8 の矢印がつながる。}
\label{fig:rlb-chain}
\end{figure}

この章の方法と、学習するものとの関係も、ここで先に述べておきます。
この章の後半のコードが実際に書き換えるのは、方策 $\pi$ ではなく、「状態や行動がどれだけ良いか」を表す数の表です（\secref{sec:rlb-value}の価値関数 $V$ と $Q$）。
方策は、その表から「いちばん良い行動を選ぶ」ことで作ります。
こうして作った最適方策は、各マスで 1 つの向きを確率 1 で選ぶ決定的方策です（\figref{fig:rlb-chain}の右）。
その場合、軌跡に残る偶然は、環境の $P$ だけになります。
方策を確率のまま、パラメータを持つ関数で表して直接学ぶ方法は、次章（\chapref{ch:policy-gradient}）で扱います。
LLM の強化学習は、こちらの方法を使います。

最後に、もう一方の確率モデル $P$ について補足します。
この章では、$P$ を学習しません。
しかし、実際に動いた経験から $P$（と報酬）を推定し、推定した $P$ を使って方策を考える方法もあります。
これを\term{モデルベース強化学習}{model-based reinforcement learning}と呼びます（「モデル」は環境の $P$ と $r$ のこと。\secref{sec:rlb-mc-td}）。
本書では扱いません。

この節で分かったことをまとめます。
エージェントの振る舞いは方策 $\pi(a\mid s)$ で、1 回分の経験は軌跡 $\tau$ で表します。
良さの基準は、割引した報酬の合計である収益 $G_t$ で、$G_t=r_t+\gamma G_{t+1}$ という再帰式が成り立ちます。
目的は、期待収益 $J(\pi)=\E[G_0]$ を最大にすることです。
確率モデルは 2 つあり、環境の遷移確率 $P$ は与えられたまま変わらず、学習で良くするのは方策 $\pi$ です。

\section{価値関数とベルマン方程式}\label{sec:rlb-value}

前節の最後に、方策の良さは収益の期待値で測ると述べました。
この節の問いは、「収益の期待値を状態ごとに表すにはどうするか。そして、それをどう計算するか」です。

\subsection*{状態の価値と、行動の価値}

格子世界で、ゴールの隣のマスにいるのと、穴の隣のマスにいるのとでは、前者のほうが「良い」はずです。
この良さを、数で表します。

\begin{definition}[価値関数]\label{def:rlb-value}
方策 $\pi$ の\term{状態価値関数}{state-value function} $V^\pi$ と\term{行動価値関数}{action-value function} $Q^\pi$ を、次の式で定義します。
\begin{align}
V^\pi(s)&\defeq\E_\pi[\,G_t\mid s_t=s\,] \label{eq:rlb-v-def}\\
Q^\pi(s,a)&\defeq\E_\pi[\,G_t\mid s_t=s,\ a_t=a\,] \label{eq:rlb-q-def}
\end{align}
\end{definition}

記号を説明します。
$\E_\pi[\,\cdot\mid s_t=s\,]$ は、条件付きの期待値です。
条件は、「状態 $s$ から出発して、そのあとずっと方策 $\pi$ で行動する」です。
この条件のもとで起こりうるすべての軌跡について、収益を、確率を重みにして平均します。
つまり $V^\pi(s)$ は、「状態 $s$ から方策 $\pi$ で動いたときに、平均してどれだけの収益が得られるか」です。
$Q^\pi(s,a)$ は、「状態 $s$ で、最初の 1 回だけは行動 $a$ を選び、そのあとは方策 $\pi$ で動いたときの、平均の収益」です。
左辺に時刻 $t$ がないのには、理由があります。
マルコフ性があり、環境のルールも方策も時刻で変わりません。
そのため、何時に状態 $s$ にいても、そこから先の平均は同じになります。
終端状態ではもう報酬をもらえないので、終端状態の価値は 0 と決めます。

前節の実験では、ランダム方策でスタート（状態 8）から 1 万回動かした収益の平均が $-0.066$ でした。
これは、ランダム方策の $V^\pi(8)$ を、実験で見積もった値です。
スタートは 8 と決まっているので、$V^\pi(8)$ は前節の期待収益 $J(\pi)$ そのものです。

2 つの価値関数は、次の式で結ばれています。
\begin{equation}
V^\pi(s)=\sum_{a}\pi(a\mid s)\,Q^\pi(s,a)
\label{eq:rlb-v-from-q}
\end{equation}
状態 $s$ では、確率 $\pi(a\mid s)$ で行動 $a$ が選ばれます。
行動 $a$ が選ばれたあとの平均の収益が $Q^\pi(s,a)$ です。
それを、選ばれる確率を重みにして平均すれば、$V^\pi(s)$ になります。

\subsection*{ベルマン期待方程式を導く}

定義のままでは、$V^\pi$ を計算できません。
起こりうる軌跡は無数にあり、すべてを書き出すことはできないからです。
そこで、収益の再帰式 $G_t=r_t+\gamma G_{t+1}$ を使って、「隣どうしの状態の価値の関係」を作ります。
考え方は、\figref{fig:rlb-backup}のとおりです。
状態 $s$ からの 1 歩だけを場合分けし、そこから先は「次の状態の価値」にまとめます。

\begin{figure}[tbp]
\centering
\begin{tikzpicture}[
  st/.style={draw=bookblue,fill=figblue,circle,minimum size=7mm,inner sep=0pt,font=\small},
  ac/.style={fill=black!70,circle,minimum size=2.4mm,inner sep=0pt}]
  \node[st] (s) at (0,0) {$s$};
  \node[ac,label={[lbl]left:$a_1$}] (a1) at (-2.0,-1.3) {};
  \node[ac,label={[lbl]right:$a_2$}] (a2) at (2.0,-1.3) {};
  \node[st] (s1) at (-3.0,-2.8) {$s'_1$};
  \node[st] (s2) at (-1.0,-2.8) {$s'_2$};
  \node[st] (s3) at (1.0,-2.8) {$s'_3$};
  \node[st] (s4) at (3.0,-2.8) {$s'_4$};
  \foreach \i in {1,2,3,4} \node[lbl] at ($(s\i)+(0,-0.6)$) {$V^\pi(s'_\i)$};
  \draw[thick,black!70] (s) -- node[above left,lbl]{$\pi(a\mid s)$} (a1);
  \draw[thick,black!70] (s) -- (a2);
  \draw[thick,black!70] (a1) -- node[left=2.5mm,lbl,align=right]{確率 $P(s'\mid s,a)$\\報酬 $r$} (s1);
  \draw[thick,black!70] (a1) -- (s2);
  \draw[thick,black!70] (a2) -- (s3);
  \draw[thick,black!70] (a2) -- (s4);
  \node[lbl,anchor=west] at (3.5,0) {状態 $s$ の価値 $V^\pi(s)$ を知りたい};
  \node[lbl,anchor=west] at (3.5,-1.3) {方策が行動 $a$ を選ぶ};
  \node[lbl,anchor=west,align=left] at (3.5,-2.8) {環境が次の状態 $s'$ と報酬 $r$ を返す\\ここから先は $\gamma V^\pi(s')$ にまとめる};
\end{tikzpicture}
\caption{ベルマン期待方程式の考え方（模式図）。丸が状態、黒い点が行動。
行動を 2 つ（$a_1,a_2$）、それぞれの次の状態を 2 つだけ描いた（格子世界では行動は 4 つ、次の状態は最大 3 つ）。
確率と報酬のラベルは左端の枝にだけ書いたが、どの枝にも付く。
1 歩先の場合分けをすべて並べ、それぞれの「報酬 $+$ $\gamma\times$次の状態の価値」を、確率 $\pi(a\mid s)P(s'\mid s,a)$ で平均する。}
\label{fig:rlb-backup}
\end{figure}

1 行ずつ導きます。
\begin{align}
V^\pi(s)
&=\E_\pi[\,G_t\mid s_t=s\,] \notag\\
&=\E_\pi[\,r_t+\gamma G_{t+1}\mid s_t=s\,] \notag\\
&=\sum_a\pi(a\mid s)\sum_{s'}P(s'\mid s,a)\;
  \E_\pi[\,r_t+\gamma G_{t+1}\mid s_t=s,a_t=a,s_{t+1}=s'\,] \notag\\
&=\sum_a\pi(a\mid s)\sum_{s'}P(s'\mid s,a)\,
  \bigl[\,r(s,a,s')+\gamma\,\E_\pi[\,G_{t+1}\mid s_{t+1}=s'\,]\,\bigr] \notag\\
&=\sum_a\pi(a\mid s)\sum_{s'}P(s'\mid s,a)\,
  \bigl[\,r(s,a,s')+\gamma\,V^\pi(s')\,\bigr]
\label{eq:rlb-bellman-v}
\end{align}
各行の理由は、次のとおりです。
\begin{itemize}
\item 1 行目は、$V^\pi$ の定義（式\eqref{eq:rlb-v-def}）です。
\item 2 行目は、収益の再帰式（式\eqref{eq:rlb-return-recursive}）を代入しました。
\item 3 行目は、最初の 1 歩で場合分けしました。
行動 $a$ が選ばれ、次の状態が $s'$ になる確率は $\pi(a\mid s)\,P(s'\mid s,a)$ です。
全体の平均は、「場合ごとの平均」を、その場合が起こる確率で重み付けして足せば求まります。
\item 4 行目では、まず、期待値の線形性（式\eqref{eq:math-linearity}）を使って式を分けました。
和の平均は平均の和になり、定数倍（ここでは $\gamma$）は平均の外に出せます。
$s$、$a$、$s'$ が決まれば、報酬 $r_t$ は $r(s,a,s')$ という決まった数です。
また、マルコフ性により、$s'$ から先の収益は $s'$ だけで決まり、その前の $s$ や $a$ にはよりません。
\item 5 行目は、$V^\pi$ の定義をもう一度使いました。$\E_\pi[G_{t+1}\mid s_{t+1}=s']$ は、$V^\pi(s')$ そのものです。
\end{itemize}

式\eqref{eq:rlb-bellman-v}を、\term{ベルマン期待方程式}{Bellman expectation equation}と呼びます\cite{bellman1957dynamic}。
言葉にすると、こうなります。
\begin{center}
いまの状態の価値 $=$（すぐにもらう報酬 $+$ $\gamma\times$次の状態の価値）の平均
\end{center}

行動価値についても、同じ考え方で式が作れます。
行動 $a$ はもう決まっているので、平均を取るのは次の状態 $s'$ についてだけです。
\begin{equation}
Q^\pi(s,a)=\sum_{s'}P(s'\mid s,a)\,\bigl[\,r(s,a,s')+\gamma\,V^\pi(s')\,\bigr]
\label{eq:rlb-q-from-v}
\end{equation}
式\eqref{eq:rlb-q-from-v}を式\eqref{eq:rlb-v-from-q}に代入すると、式\eqref{eq:rlb-bellman-v}に戻ります。

\subsection*{数値で確かめる}

式\eqref{eq:rlb-q-from-v}を、格子世界の数で確かめます。
状態 2（ゴールの左隣）で「右」を選ぶと、結果は 3 通りです。
確率 0.8 でゴールに入り、報酬 $+1$ をもらって終わります（終端状態の価値は 0）。
確率 0.1 で上にそれて、外枠にぶつかり、状態 2 にとどまります。
確率 0.1 で下にそれて、状態 6 に進みます。
ランダム方策の価値 $V^\pi(2)=0.2616$、$V^\pi(6)=-0.3371$ は、すぐあとの\figref{fig:rlb-v-random}で求める値を、先に借ります。
\begin{align*}
Q^\pi(2,\text{右})
&=0.8\,(1+0.9\times0)+0.1\,(0+0.9\times0.2616)+0.1\,(0+0.9\times(-0.3371))\\
&=0.8+0.0235-0.0303=0.7932
\end{align*}
ほかの 3 つの行動も同じように計算すると、「上」が $0.2998$、「下」が $-0.1313$、「左」が $0.0848$ です。
式\eqref{eq:rlb-v-from-q}のとおりに確率 $1/4$ で平均すると、$(0.2998+0.7932-0.1313+0.0848)/4=0.2616$ で、$V^\pi(2)$ と一致します。

\subsection*{方策評価：方程式を繰り返しで解く}

ここから\secref{sec:rlb-dp}の終わりまでは、特別に、遷移確率の表 $P$ を見てよいことにします。
まず「答えを計算で出せる場合」の解き方を知り、\secref{sec:rlb-mc-td}から、表を見ずに同じ答えに近づく方法へ進みます。

ベルマン期待方程式は、9 個の状態の価値 $V^\pi(s)$ を未知数とする、9 本の連立方程式です。
これを解いて $V^\pi$ を求めることを、\term{方策評価}{policy evaluation}と呼びます。
9 本なら、連立一次方程式として直接解くこともできます。
ただ、状態が多いと直接解く計算は重くなります。
また、次節の価値反復は「平均」が「最大」に変わるので、連立一次方程式ではなくなります。
そこで、どちらにも使える繰り返しの解き方を使います。
価値の表 $V$ を、すべて 0 で始めます。
式\eqref{eq:rlb-bellman-v}の右辺にいまの表 $V$ を入れて計算し、その結果を新しい表にします。
これを、表がほとんど変わらなくなるまで繰り返します。
この方法を\jterm{反復方策評価}と呼びます。
繰り返すたびに、表は正しい $V^\pi$ に近づきます（理由は\secref{sec:rlb-dp}で述べます）。

\snippet{反復方策評価。ベルマン期待方程式の右辺を、表の書き換えの規則として使う}{code:rlb-policy-eval}{ch20/policy_eval.py}{policy-eval}

\code{q_from_v} は、式\eqref{eq:rlb-q-from-v}をそのまま書いたものです。
\code{P[s][a]} から（確率, 次の状態, 報酬）を 1 組ずつ取り出し、「確率 $\times$（報酬 $+$ $\gamma\times$次の状態の価値）」を足し合わせます。
12〜13 行目が、式\eqref{eq:rlb-v-from-q}です。
方策の確率 \code{policy[s][a]} を重みにして、4 つの行動の価値を平均します。
14 行目で、表の変化の最大値を測り、16 行目で、それが十分小さければ繰り返しを終えます。
「十分小さい」の基準は、引数 \code{tol} の $10^{-8}$ です。

ランダム方策で、最初の書き換えを手で追ってみます。
表はすべて 0 なので、残るのは報酬の平均だけです。
状態 2 で「右」を選ぶと、確率 0.8 でゴールに入るので、報酬の平均は $0.8$ です。
「上」と「下」では、確率 0.1 でゴールの側にそれるので $0.1$、「左」では 0 です。
4 つを確率 $1/4$ で平均すると、$(0.1+0.8+0.1+0)/4=0.25$ になります。
穴の隣の状態 6 と 11 は、同じ計算で $-0.25$ です。
残りの状態は、1 回目では 0 のままです。
2 回目には、状態 2 の $0.25$ が左隣の状態 1 に伝わり、$0.056$ が入ります。
こうして、書き換えのたびに、値が 1 マスずつ遠くへ伝わっていきます（ここでの「伝わる」は、そのマスの値が 0 から動くことです）。

ランダム方策を評価すると、98 回の書き換えで収束しました。
結果が\figref{fig:rlb-v-random}です。

\begin{figure}[tbp]
\centering
\begin{tikzpicture}[x=1.3cm,y=0.86cm]
  \rlbcell{0}{0}{+0.05}{bookgreen!3}
\rlbcell{1}{0}{+0.13}{bookgreen!9}
\rlbcell{2}{0}{+0.26}{bookgreen!18}
\rlbgoal{3}{0}
\rlbcell{0}{1}{-0.01}{bookaccent!0}
\rlbwall{1}{1}
\rlbcell{2}{1}{-0.34}{bookaccent!24}
\rlbpit{3}{1}
\rlbcell{0}{2}{-0.07}{bookaccent!5}
\rlbcell{1}{2}{-0.15}{bookaccent!11}
\rlbcell{2}{2}{-0.31}{bookaccent!22}
\rlbcell{3}{2}{-0.58}{bookaccent!41}

  \rlbframe
  \rlbnums
\end{tikzpicture}
\caption{ランダム方策の状態価値 $V^\pi$（$\gamma=0.9$）。緑が正、赤が負で、色が濃いほど絶対値が大きい。
G はゴール、灰色は壁。ゴールに近い状態は価値が高く、穴に近い状態は低い。}
\label{fig:rlb-v-random}
\end{figure}

ゴールの左隣（状態 2）の価値は $+0.26$、穴の下（状態 11）は $-0.58$ です。
でたらめに歩くと、穴の近くでは落ちる見込みが高いことが、数に表れています。
スタート（状態 8）の価値は $-0.066$ です。
前節で 1 万エピソードの収益を平均した値 $-0.066$ と、小数第 3 位まで一致します。
実験をしなくても、遷移確率の表さえあれば、計算で同じ値が求まるわけです。

この節で分かったことをまとめます。
状態価値 $V^\pi(s)$ は、状態 $s$ から方策 $\pi$ で動いたときの収益の期待値です。
ベルマン期待方程式は、いまの状態の価値を、1 歩先の状態の価値で表します。
この式を繰り返し使うと、$V^\pi$ を計算できます。

\section{最適な方策を求める：動的計画法}\label{sec:rlb-dp}

前節では、与えられた方策の価値を計算しました。
しかし、本当に知りたいのは、いちばん良い方策です。
この節の問いは、「最良の方策とは何か。遷移確率が分かっているとき、それをどう求めるか」です。
遷移確率を使い、ベルマン方程式で価値の表を順に埋めていく方法を、まとめて\term{動的計画法}{dynamic programming}と呼びます\cite{bellman1957dynamic}。
前節の反復方策評価も、その 1 つです。

\subsection*{最適価値関数とベルマン最適方程式}

あらゆる方策の中で達成できる最大の価値を、\jterm{最適価値関数}と呼びます。
\begin{equation}
V^\ast(s)\defeq\max_\pi V^\pi(s),\qquad Q^\ast(s,a)\defeq\max_\pi Q^\pi(s,a)
\label{eq:rlb-optimal}
\end{equation}
すべての状態で $V^\pi(s)=V^\ast(s)$ となる方策を、\term{最適方策}{optimal policy}と呼び、$\pi^\ast$ と書きます。
状態と行動が有限個の MDP には（割引率 $\gamma<1$ のとき。この章の設定です）、最適方策が必ずあり、決定的方策の中に見つかることが知られています\cite{sutton2018rl}。
最適方策は、すべての状態で同時に価値を最大にするので、とくにスタートでの価値、つまり期待収益 $J$（式\eqref{eq:rlb-objective}）も最大にします。
\secref{sec:rlb-return}で決めた目的と、同じものを指しているわけです。

ここからの議論の骨組みを、先に示します。
(1) 最適方策があることは、上に書いた、知られている事実です（証明は省きます）。
(2) 最適方策は、各状態で、いちばん価値の高い行動を選びます。
(3) そのため、期待方程式の「方策の確率で平均する」が「最大を取る」に変わります。
(4) こうしてできる式が、ベルマン最適方程式です。

(2) の理由を述べます。
最適方策は、各状態で、いちばん価値の高い行動を選びます。
価値の低い行動に確率を割いても、得をしないからです。
そこで、式\eqref{eq:rlb-v-from-q}の「方策の確率で平均する」を「最大のものを選ぶ」に置き換えると、次の 1 つ目の式になります。
2 つ目は、式\eqref{eq:rlb-q-from-v}と同じ考え方で作った式です。
\[
V^\ast(s)=\max_a Q^\ast(s,a),\qquad
Q^\ast(s,a)=\sum_{s'}P(s'\mid s,a)\,\bigl[\,r(s,a,s')+\gamma\,V^\ast(s')\,\bigr]
\]
2 つ目を 1 つ目に代入すると式\eqref{eq:rlb-bellman-opt-v}、1 つ目を 2 つ目に代入すると式\eqref{eq:rlb-bellman-opt-q}になります。
これが\term{ベルマン最適方程式}{Bellman optimality equation}です。
$a'$ は、次の状態 $s'$ で選ぶ行動を表します。
\begin{align}
V^\ast(s)&=\max_a\sum_{s'}P(s'\mid s,a)\,\bigl[\,r(s,a,s')+\gamma\,V^\ast(s')\,\bigr]
\label{eq:rlb-bellman-opt-v}\\
Q^\ast(s,a)&=\sum_{s'}P(s'\mid s,a)\,\bigl[\,r(s,a,s')+\gamma\max_{a'}Q^\ast(s',a')\,\bigr]
\label{eq:rlb-bellman-opt-q}
\end{align}
\tabref{tab:rlb-bellman}に、期待方程式と並べます。
違いは、行動についての「平均」が「最大」に変わったことだけです。

\begin{table}[tbp]
\centering
\caption{2 種類のベルマン方程式。$[\cdots]$ は $r(s,a,s')+\gamma\times$（次の状態の価値）を表す。}
\label{tab:rlb-bellman}
\small
\begin{tabular}{lll}
\toprule
 & 期待方程式（方策 $\pi$ を決めたとき） & 最適方程式（最良の方策） \\
\midrule
状態価値 & $V^\pi(s)=\sum_a\pi(a\mid s)\sum_{s'}P\,[\cdots]$ & $V^\ast(s)=\max_a\sum_{s'}P\,[\cdots]$ \\
行動についての操作 & 方策の確率で平均する & 最大のものを選ぶ \\
解き方 & 反復方策評価 & 価値反復 \\
\bottomrule
\end{tabular}
\end{table}

$Q^\ast$ が分かれば、最適方策はすぐに作れます。
各状態で、$Q^\ast$ が最大の行動を選ぶだけです。
\begin{equation}
\pi^\ast(s)=\argmax_a Q^\ast(s,a)
\label{eq:rlb-greedy}
\end{equation}
このように、価値がいちばん高い行動を必ず選ぶ方策を、その価値に対する\term{貪欲方策}{greedy policy}と呼びます。

\subsection*{方策改善と方策反復}
\skimread{押さえる結論は、「方策の評価と改善を交互に繰り返すと、最適方策に着く。格子世界では 2 回の改善で着く」です。途中の数字は 2 周目で十分です。}

\figref{fig:rlb-v-random}のランダム方策を、良くすることを考えます。
状態 2 では、4 つの行動のうち「右」の価値 $Q^\pi(2,\text{右})=0.79$ が最大でした。
それなら、状態 2 では必ず「右」を選ぶように変えれば、方策は良くなるはずです。
すべての状態で同じように、$Q^\pi$ が最大の行動を選ぶ方策に取り替えることを、\term{方策改善}{policy improvement}と呼びます。
改善した方策は、どの状態でも元の方策より悪くならないことが証明されています\cite{sutton2018rl}。

方策評価と方策改善を交互に繰り返す方法が、\term{方策反復}{policy iteration}です（\figref{fig:rlb-gpi}）。
方策が変わらなくなったら、それが最適方策です。
ランダム方策 $\pi_0$ から始めた結果が、次の出力です（\file{ch20/value_iteration.py} の \code{policy_iteration}）。

\begin{figure}[tbp]
\centering
\begin{tikzpicture}
  \node[box,minimum width=24mm,minimum height=9mm] (pi) at (0,0) {方策 $\pi$};
  \node[greenbox,minimum width=24mm,minimum height=9mm] (v) at (7,0) {価値 $V^\pi$};
  \draw[arr] (pi.north east) to[bend left=14] node[above,font=\small]{方策評価：$\pi$ の価値を計算する} (v.north west);
  \draw[arr] (v.south west) to[bend left=14] node[below,font=\small]{方策改善：価値が最大の行動を選ぶ方策に取り替える} (pi.south east);
\end{tikzpicture}
\caption{方策反復。評価と改善を交互に行うと、方策は悪くならずに良くなっていき、最適方策で止まる。}
\label{fig:rlb-gpi}
\end{figure}

\outputfile{ch20/out/policy_iteration.txt}

各行は、方策 $\pi_k$ の価値と、それを改善して $\pi_{k+1}$ にするときに行動が変わる状態の数です。
ランダム方策 $\pi_0$ は、どの状態でも 4 つの行動を $1/4$ ずつ選ぶので、選ぶ行動が 1 つに決まっていません。
改善すると、9 状態のすべてが「1 つの行動を確率 1 で選ぶ」に変わるので、9 個と数えています。
スタートの価値 $V(8)$ は、1 回の改善で $-0.066$ から $+0.545$ に上がりました。
2 回目の改善では状態 10 の行動だけが「左」から「上」に変わり、3 回目には何も変わりません。
たった 2 回の改善で、最適方策に着きました。

なぜ、状態 10 の行動は 1 回で決まらなかったのでしょうか。
\figref{fig:rlb-pi-improve}に、1 回目の改善でできた方策 $\pi_1$ と、2 回目の改善でできた $\pi_2$ を並べました。
1 回目の改善は、ランダム方策の価値（\figref{fig:rlb-v-random}）を見て行動を選びます。
でたらめに歩くかぎり、穴の隣の状態 6 は危険な場所で、価値は $-0.337$ です。
左隣の状態 9 の $-0.155$ より低いので、状態 10 では、「上」（行動価値 $-0.309$）より「左」（$-0.170$）が選ばれます。

ところが $\pi_1$ は、状態 6 に着いたら必ず「上」を選びます。
この方策を評価し直すと、状態 6 の価値は $0.635$ に上がり、状態 9 の $0.479$ を上回ります。
すると状態 10 でも、「上」（$0.522$）が「左」（$0.442$）を上回り、2 回目の改善で「上」に変わります。
行き先での振る舞いが良くなると、そこへ向かう行動の価値も上がるわけです。
評価と改善を交互に繰り返すのは、このためです。

\begin{figure}[tbp]
\centering
\begin{tikzpicture}[x=1.05cm,y=0.72cm]
  \rlbarrow{0}{0}{0}
\rlbarrow{1}{0}{0}
\rlbarrow{2}{0}{0}
\rlbgoal{3}{0}
\rlbarrow{0}{1}{90}
\rlbwall{1}{1}
\rlbarrow{2}{1}{90}
\rlbpit{3}{1}
\rlbarrow{0}{2}{90}
\rlbarrow{1}{2}{180}
\rlbarrow{2}{2}{180}
\rlbarrow{3}{2}{180}

  \rlbframe
  \rlbnums
  \draw[bookaccent,very thick] (2,-2) rectangle ++(1,-1);
  \node[font=\small] at (2,0.4) {方策 $\pi_1$（改善 1 回）};
  \begin{scope}[shift={(5.2,0)}]
    \rlbarrow{0}{0}{0}
\rlbarrow{1}{0}{0}
\rlbarrow{2}{0}{0}
\rlbgoal{3}{0}
\rlbarrow{0}{1}{90}
\rlbwall{1}{1}
\rlbarrow{2}{1}{90}
\rlbpit{3}{1}
\rlbarrow{0}{2}{90}
\rlbarrow{1}{2}{180}
\rlbarrow{2}{2}{90}
\rlbarrow{3}{2}{180}

    \rlbframe
    \rlbnums
    \draw[bookaccent,very thick] (2,-2) rectangle ++(1,-1);
    \node[font=\small] at (2,0.4) {方策 $\pi_2$（改善 2 回。最適方策）};
  \end{scope}
\end{tikzpicture}
\caption{方策反復の途中の方策。違うのは、枠で囲んだ状態 10 だけ。
$\pi_1$ は、穴の隣の状態 6 を避けて左へ向かう。$\pi_1$ の価値で見直した $\pi_2$ は、上へ向かう。}
\label{fig:rlb-pi-improve}
\end{figure}

\subsection*{価値反復}

方策反復は、評価と改善を交互に行うので、改善のたびに方策評価をやり直します。
実は、評価が収束するまで待つ必要はありません。
ベルマン最適方程式\eqref{eq:rlb-bellman-opt-v}の右辺を、表の書き換えの規則としてそのまま使えば、$V^\ast$ が直接求まります。
この方法を\term{価値反復}{value iteration}と呼びます。

\snippet{価値反復}{code:rlb-value-iteration}{ch20/value_iteration.py}{vi}

\coderef{code:rlb-policy-eval}との違いは、8 行目だけです。
方策の確率で平均していたところが、\code{max} に変わりました。
収束した $V$ から式\eqref{eq:rlb-q-from-v}で各行動の価値を計算し、最大の行動を選べば、最適方策になります（同じファイルの \code{greedy_policy}）。

実行すると、29 回の書き換えで収束しました。
方策反復では、3 回の方策評価に、それぞれ 98 回、31 回、29 回、合わせて 158 回の書き換えを使いました（\file{ch20/fig_dp_extra.py} で数えました）。
1 回の書き換えでする計算は、どちらも「各状態で 4 つの行動の価値を求める」で、ほぼ同じです。
価値反復は、方策反復のおよそ 5 分の 1 の手間で、同じ最適方策に着いたことになります。

\figref{fig:rlb-vi-sweeps}は、途中の表です。
1 回目では、ゴールの隣の状態 2 にだけ価値が入ります。
「右」を選ぶと確率 0.8 でゴールに入るので、値は $0.8\times1=0.80$ です。
穴の隣の状態 6 と 11 は、0 のままです。
穴に向かわない行動を選べば落ちずに済むので、最大を取ると、悪い結果は表に出てきません。
平均を取る方策評価では、状態 6 は 1 回目に $-0.25$ になりました（収束した値は\figref{fig:rlb-v-random}）。
平均と最大の違いが、ここに現れています。
2 回目には、左隣の状態 1 に $0.8\times0.9\times0.80=0.576$ が入ります（図では $0.58$）。
状態 6 にも、「上」を選んで状態 2 に進む分の値 $0.48$ が入ります。
図の赤い枠をコマごとに追ってください。
初めて値が入るマスは、1 回目が状態 2、2 回目が 1 と 6、3 回目が 0 と 10、4 回目が 4 と 9 と 11、5 回目が 8 です。
書き換えのたびに、値の入ったマスが 1 マスずつゴールから離れ、5 回目にスタートまで届きます。
ゴールで受け取る報酬の効果を、手前の状態の価値へ 1 マスずつ反映させていくこの仕組みが、\secref{sec:rlb-what}の信用割当の問題への、1 つの答えです。

\begin{figure}[tbp]
\centering
\begin{tikzpicture}[x=0.98cm,y=0.68cm]
  \renewcommand{\rlbfont}{\scriptsize}
  \foreach \k/\x/\y in {1/0/0,2/4.4/0,3/8.8/0,4/0/-4.0,5/4.4/-4.0} {
    \begin{scope}[shift={(\x,\y)}]
      \input{code/ch20/out/fig_vi_frame\k.tex}
      \rlbframe
      \rlbnumstag
      \node[font=\footnotesize] at (2,0.42) {\k 回目};
    \end{scope}}
  \begin{scope}[shift={(8.8,-4.0)}]
    \draw[bookaccent,very thick] (0.1,-0.9) rectangle ++(0.7,-0.7);
    \node[font=\footnotesize,anchor=west,align=left] at (1.0,-1.25) {その回に初めて\\値が入ったマス};
  \end{scope}
\end{tikzpicture}
\caption{価値反復の途中経過（1〜5 回目。\file{ch20/fig_dp_extra.py} で、\coderef{code:rlb-value-iteration}の途中の表を書き出した）。
赤い枠は、その回に初めて 0 でない値が入ったマス。
価値は、1 回の書き換えで 1 マスずつ、ゴールから遠くへ伝わる。5 回目に、スタート（状態 8）まで届く。}
\label{fig:rlb-vi-sweeps}
\end{figure}

収束した $V^\ast$ と最適方策が、\figref{fig:rlb-v-star}です。
スタートの価値は $V^\ast(8)=0.545$ で、方策反復の結果と一致します。

\begin{figure}[tbp]
\centering
\begin{tikzpicture}[x=1.25cm,y=0.85cm]
  \rlbcell{0}{0}{+0.72}{bookgreen!50}
\rlbcell{1}{0}{+0.83}{bookgreen!58}
\rlbcell{2}{0}{+0.94}{bookgreen!66}
\rlbgoal{3}{0}
\rlbcell{0}{1}{+0.63}{bookgreen!44}
\rlbwall{1}{1}
\rlbcell{2}{1}{+0.64}{bookgreen!44}
\rlbpit{3}{1}
\rlbcell{0}{2}{+0.55}{bookgreen!38}
\rlbcell{1}{2}{+0.48}{bookgreen!34}
\rlbcell{2}{2}{+0.53}{bookgreen!37}
\rlbcell{3}{2}{+0.31}{bookgreen!22}

  \rlbframe
  \rlbnums
  \node[font=\small] at (2,0.35) {最適価値関数 $V^\ast$};
  \begin{scope}[shift={(4.6,0)}]
    \rlbarrow{0}{0}{0}
\rlbarrow{1}{0}{0}
\rlbarrow{2}{0}{0}
\rlbgoal{3}{0}
\rlbarrow{0}{1}{90}
\rlbwall{1}{1}
\rlbarrow{2}{1}{90}
\rlbpit{3}{1}
\rlbarrow{0}{2}{90}
\rlbarrow{1}{2}{180}
\rlbarrow{2}{2}{90}
\rlbarrow{3}{2}{180}

    \rlbframe
    \rlbnums
    \node[font=\small] at (2,0.35) {最適方策 $\pi^\ast$};
  \end{scope}
\end{tikzpicture}
\caption{価値反復で求めた最適価値関数（左）と最適方策（右）。
スタート（左下）からは、穴から遠い上の道（$8\to4\to0\to1\to2$）を選ぶ。穴の下の状態では、穴から離れる向きを選ぶ。}
\label{fig:rlb-v-star}
\end{figure}

最適方策には、面白い点があります。
スタートからゴールへは、上の道（$8\to4\to0\to1\to2$）と下の道（$8\to9\to10\to6\to2$）があり、どちらも最短で 5 歩です。
しかし、最適方策は上の道を選びます（$Q^\ast(8,\text{上})=0.545$、$Q^\ast(8,\text{右})=0.450$。大きいほうが $V^\ast(8)$ です）。
下の道は穴の隣を通るので、滑って落ちる危険があるからです。
「床が滑る」という環境の性質が、価値の数にきちんと反映されています。

この 2 つの数を、ベルマン最適方程式で確かめます。
\figref{fig:rlb-v-star}の左の表から、状態 4、8、9 の値を使います。
小数第 4 位まで書くと、$V^\ast(4)=0.6292$、$V^\ast(8)=0.5452$、$V^\ast(9)=0.4787$ です。
状態 8 で「上」を選ぶと、確率 0.8 で状態 4 に進み、確率 0.1 ずつで、状態 8 にとどまるか、右の状態 9 にそれます。
報酬はどれも 0 なので、残るのは「確率 $\times\,\gamma\times$次の状態の価値」の和です（「右」も同様）。
\begin{align*}
Q^\ast(8,\text{上})&=0.8\times0.9\times0.6292+0.1\times0.9\times0.5452+0.1\times0.9\times0.4787=0.5452\\
Q^\ast(8,\text{右})&=0.8\times0.9\times0.4787+0.1\times0.9\times0.6292+0.1\times0.9\times0.5452=0.4504
\end{align*}
残りの 2 つは、「下」が $0.4847$、「左」が $0.4982$ です。
4 つの中の最大は「上」の $0.5452$ で、$V^\ast(8)$ と一致します（式\eqref{eq:rlb-bellman-opt-v}のとおりです）。

\subsection*{床の滑りやすさを変えると、最適方策はどう変わるか}
\skimread{押さえる結論は、「滑りやすいほど、最適方策は遠回りでも穴に落ちない行動を選び、スタートの価値 $V^\ast(8)$ は下がる」です。}

上の道が選ばれた理由は、床が滑ることでした。
では、滑りやすさを変えると、最適方策はどう変わるでしょうか。
\file{ch20/sweep_slip.py} で、左右にそれる確率（\code{SLIP}。これまでは 0.1）を 0、0.1、0.2、0.3 の 4 通りに変え、そのたびに遷移確率の表を作り直して価値反復を走らせました（$\gamma=0.9$ は固定）。
結果が\figref{fig:rlb-slip}です。
上の段の 4 枚の矢印を左から順に見比べ、下のグラフで $V^\ast(8)$ の下がり方を見てください。

\begin{figure}[tbp]
\centering
\begin{tikzpicture}[x=0.72cm,y=0.62cm]
  \foreach \k/\slp in {0/0,1/0.1,2/0.2,3/0.3} {
    \begin{scope}[shift={(4.45*\k,0)}]
      \input{code/ch20/out/fig_slip\k_pi.tex}
      \rlbframe
      \rlbnums
      \node[font=\footnotesize] at (2,0.45) {SLIP $=\slp$};
    \end{scope}}
\end{tikzpicture}

\medskip
\begin{tikzpicture}
\begin{axis}[height=0.3\linewidth,width=0.62\linewidth,xlabel={左右にそれる確率 SLIP},ylabel={スタートの価値},
  xmin=-0.02,xmax=0.32,ymin=0,ymax=0.75,xtick={0,0.1,0.2,0.3},
  legend pos=outer north east,legend cell align=left]
\addplot[bookblue,mark=*] table[col sep=comma,x=slip,y=q8up]{code/ch20/out/sweep_slip.csv};
\addplot[bookaccent,mark=square*,dashed] table[col sep=comma,x=slip,y=q8right]{code/ch20/out/sweep_slip.csv};
\legend{{$Q^\ast(8,\text{上})$（上の道）},{$Q^\ast(8,\text{右})$（下の道）}}
\end{axis}
\end{tikzpicture}
\caption{床の滑りやすさ SLIP を 4 通りに変えたときの最適方策（上の段）と、スタートで「上」と「右」を選んだときの価値（下。大きいほうが $V^\ast(8)$）。$\gamma=0.9$ は固定（\file{ch20/sweep_slip.py}）。
SLIP $=0$ では 2 つの価値が等しく、上の段の状態 8 の矢印は、同点のうち先に数えた「上」を描いた。}
\label{fig:rlb-slip}
\end{figure}

滑らない床（SLIP $=0$）では、上の道も下の道も 5 歩で確実にゴールに着くので、2 つの価値は $0.9^4=0.6561$ で等しくなります。
滑りやすくなるほど、最適方策は、遠回りになっても穴に落ちない行動を選ぶようになります。
SLIP $=0.2$ では、状態 6 で、壁に向かう「左」を選びます。
壁にぶつかればその場にとどまり、それる先は上（状態 2）か下（状態 10）なので、穴には落ちません。
状態 11 でも、外枠に向かう「下」を選んで、穴を避けます。
SLIP $=0.3$ では、ゴールの隣の状態 2 でさえ、外枠に向かう「上」を選びます。
右手にそれればゴールに入り、左手にそれても状態 1 に戻るだけで、穴の隣の状態 6 へは進みません。
実行結果（\file{ch20/out/sweep_slip.txt}）によると、最適方策で動いたときにゴールに着く確率は、SLIP $=0.1$ では 0.986、0.2 と 0.3 では 1.000 です。
安全な行動はゴールまでの歩数が増えるので、その分、割引で価値が下がります。
$V^\ast(8)$ は、0.656、0.545、0.388、0.228 と下がっていきます。
最適方策は、環境の遷移確率 $P$ によって変わります。
$P$ が変われば、同じ迷路でも「良い進み方」が変わるのです。

\subsection*{なぜ収束するのか}
\laterread{縮小写像という性質の紹介です。あとの章では使いません。収束の理由が気になったときに戻れば足ります。飛ばすときは、すぐあとのコラム「ここまでの道のり」へ進んでください。}

\figref{fig:rlb-vi-conv}は、1 回の書き換えで表がどれだけ変わったかを示します。
縦軸は対数目盛で、変化はほぼ直線的に、つまり毎回ほぼ一定の倍率で小さくなっています。

\begin{figure}[tbp]
\centering
\begin{tikzpicture}
\begin{axis}[height=0.33\linewidth,ymode=log,xlabel={書き換えの回数 $k$},ylabel={表の変化の最大値},
  xmin=0,xmax=30,legend pos=south west]
\addplot[bookblue,mark=*,mark size=1pt] table[col sep=comma,x=sweep,y=diff]{code/ch20/out/vi_convergence.csv};
\addplot[bookaccent,dashed,domain=1:29,samples=2] {0.8*0.9^(x-1)};
\legend{価値反復の実測,毎回 $0.9$ 倍になる線}
\end{axis}
\end{tikzpicture}
\caption{価値反復の収束。表の変化は、保証されている「毎回 $\gamma=0.9$ 倍以下」（破線）よりも速く小さくなる。}
\label{fig:rlb-vi-conv}
\end{figure}

これは偶然ではありません。
価値反復の 1 回の書き換えで、表と $V^\ast$ との差の最大値は、必ず $\gamma$ 倍以下になることが証明できます。
この性質を持つ操作を\term{縮小写像}{contraction mapping}と呼びます。
$\gamma=0.9$ なら、$k$ 回の書き換えで差は $0.9^k$ 倍以下になり、0 に近づきます。
\figref{fig:rlb-vi-conv}の「表の変化」が小さくなるのも、反復方策評価が収束するのも、同じ理由です。
本書では証明を省きます（強化学習の理論の専門書にあります）。

\begin{column}{ここまでの道のり}
\textbf{\secref{sec:rlb-what}〜\secref{sec:rlb-return}}：強化学習の問題を MDP として書いた。良さの基準は収益 $G_t=r_t+\gamma G_{t+1}$。\\
\textbf{\secref{sec:rlb-value}}：収益の期待値が価値関数。ベルマン期待方程式で、隣の状態の価値と結びつく。\\
\textbf{\secref{sec:rlb-dp}}：平均を最大に変えるとベルマン最適方程式。繰り返し当てはめると、最適方策が求まる。\\
\textbf{この先}：ここまでは、遷移確率 $P$ の表を使った。
この先は、$P$ を知らずに経験だけから学ぶ（\secref{sec:rlb-mc-td}、\secref{sec:rlb-q}）。
最後に、表が作れないほど大きな問題を考える（\secref{sec:rlb-fa}）。
\end{column}

この節で分かったことをまとめます。
最適価値関数は、ベルマン最適方程式を満たします。
遷移確率が分かっていれば、この式を繰り返し当てはめるだけで、最適な価値と方策が求まります。

\section{経験だけから学ぶ：モンテカルロ法と TD 学習}\label{sec:rlb-mc-td}

動的計画法は、遷移確率 $P$ と報酬関数 $r$ の表を使いました。
この 2 つを合わせて、\jterm{環境のモデル}と呼びます。
これまでの章の「モデル」（学習するネットワーク）とは別の意味なので、この章では必ず「環境のモデル」と書きます。
自転車の練習では、環境のモデルは手に入りません。
LLM の強化学習では、遷移（トークンを列の末尾に付ける）は分かっています。
しかし、状態の数が多すぎて、動的計画法のように全部の状態の表を作ることができません（\secref{sec:rlb-fa}）。
どちらの場合も、実際に動いた経験から学ぶ方法が要ります。
この節の問いは、「環境のモデルを使わずに、実際に動いて得た経験だけから価値を求めるにはどうするか」です。
環境のモデルを使わない方法を、\term{モデルフリー}{model-free}の方法と呼びます。
これに対して、環境のモデルを使う方法を、モデルベースの方法と呼びます。
モデルが与えられている場合（前節の動的計画法）のほかに、経験から $P$ と $r$ を推定して使う場合（\secref{sec:rlb-return}の最後で触れたモデルベース強化学習）があります。

\subsection*{モンテカルロ法：実際の収益を平均する}

$V^\pi(s)$ は、収益の期待値でした。
期待値は、サンプルをたくさん集めて平均すれば近似できます（\chapref{ch:math}の大数の法則）。
方策 $\pi$ で何度もエピソードを動かし、状態 $s$ を通るたびに、そこから先の収益 $G_t$ を記録して平均します。
この方法を\term{モンテカルロ法}{Monte Carlo method}と呼びます。
\secref{sec:rlb-return}で、1 万エピソードの平均から $-0.066$ を得たのが、まさにこの方法です。

「何度も動かして、通るたびに集めて平均する」を絵にしたのが、\figref{fig:rlb-mc-spacetime}の時空図です。
横軸が時刻 $t$、縦軸がマスで、1 本の折れ線が 1 回のエピソード（試行）です。
上の図では、ランダム方策の試行を 4 本重ねました（\file{ch20/fig_mc_spacetime.py}）。
見てほしいのは、状態 9 の横の帯です。
どの試行も、状態 9 を通った時刻に白い丸があります。
丸ごとに、その時刻から終わりまでの収益 $G_t$ を計算し、右の一覧に集めます。
試行 3 は、時刻 3、4、5 と続けて状態 9 にいた（2 回、壁や外枠にぶつかってとどまった）ので、3 個の値を出します。
コードの \code{mc_update} と同じく、同じエピソードの中で同じ状態を何度通っても、通るたびに 1 個のデータとして数えます。
集めた値を平均したものが、$V^\pi(9)$ の見積もりです。

4 本だけでは、平均は $+0.21$ で、正しい値 $-0.155$ から大きく外れています。
4 本のうち 2 本がたまたまゴールに着いたからです。
下の図は、試行を 2,000 本まで増やしたときの見積もりの動きです。
集めた値を全部同じ重み（$1/n$）で平均すると（青）、データが増えるにつれて揺れが小さくなり、2,000 本で $-0.157$ と正しい値に近づきます。

\begin{figure}[tbp]
\centering
\begin{tikzpicture}[x=0.52cm,y=0.5cm]
  \fill[bookorange!15] (-0.3,-8.45) rectangle (12.4,-7.55);
  \foreach \s/\r in {0/0,1/1,2/2,4/4,6/5,8/7,9/8,10/9,11/10}
    \node[font=\scriptsize,anchor=east] at (-0.35,-\r) {\s};
  \node[font=\scriptsize,anchor=east,text=bookgreen] at (-0.35,-3) {3 ゴール};
  \node[font=\scriptsize,anchor=east,text=bookaccent] at (-0.35,-6) {7 穴};
  \foreach \r in {0,...,10} \draw[black!12] (-0.3,-\r) -- (12.4,-\r);
  \draw[->,black!60] (-0.3,-10.8) -- (12.6,-10.8) node[right,font=\scriptsize]{時刻 $t$};
  \foreach \t in {0,2,...,12} \node[font=\scriptsize,anchor=north] at (\t,-10.8) {\t};
  \node[font=\scriptsize,rotate=90] at (-3.6,-5) {マス（状態）};
  \draw[bookblue,thick] (0,-7.18) -- (1,-8.18) -- (2,-9.18) -- (3,-10.18) -- (4,-6.18);
\fill[bookblue] (4,-6.18) ++(-2pt,-2pt) rectangle ++(4pt,4pt);
\draw[bookblue,fill=white,thick] (1,-8.18) circle (2.2pt);
\draw[bookorange,thick] (0,-7.06) -- (1,-8.06) -- (2,-9.06) -- (3,-10.06) -- (4,-10.06) -- (5,-10.06) -- (6,-9.06) -- (7,-10.06) -- (8,-9.06) -- (9,-10.06) -- (10,-6.06);
\fill[bookorange] (10,-6.06) ++(-2pt,-2pt) rectangle ++(4pt,4pt);
\draw[bookorange,fill=white,thick] (1,-8.06) circle (2.2pt);
\draw[bookgreen,thick] (0,-6.94) -- (1,-6.94) -- (2,-6.94) -- (3,-7.94) -- (4,-7.94) -- (5,-7.94) -- (6,-8.94) -- (7,-4.94) -- (8,-4.94) -- (9,-1.94) -- (10,-2.94);
\fill[bookgreen] (10,-2.94) ++(-2pt,-2pt) rectangle ++(4pt,4pt);
\draw[bookgreen,fill=white,thick] (3,-7.94) circle (2.2pt);
\draw[bookgreen,fill=white,thick] (4,-7.94) circle (2.2pt);
\draw[bookgreen,fill=white,thick] (5,-7.94) circle (2.2pt);
\draw[bookpurple,thick] (0,-6.82) -- (1,-7.82) -- (2,-7.82) -- (3,-7.82) -- (4,-6.82) -- (5,-6.82) -- (6,-6.82) -- (7,-3.82) -- (8,0.18) -- (9,-0.82) -- (10,-0.82) -- (11,-1.82) -- (12,-2.82);
\fill[bookpurple] (12,-2.82) ++(-2pt,-2pt) rectangle ++(4pt,4pt);
\draw[bookpurple,fill=white,thick] (1,-7.82) circle (2.2pt);
\draw[bookpurple,fill=white,thick] (2,-7.82) circle (2.2pt);
\draw[bookpurple,fill=white,thick] (3,-7.82) circle (2.2pt);

  % 終わりの印に試行の番号を添える（時刻 0〜3 で試行 1 と 2 の線が重なるため）
  \node[font=\scriptsize,text=bookblue,anchor=west] at (4.15,-6.18) {試行 1};
  \node[font=\scriptsize,text=bookorange,anchor=west] at (10.15,-6.06) {試行 2};
  \node[font=\scriptsize,text=bookgreen,anchor=north] at (10,-3.1) {試行 3};
  \node[font=\scriptsize,text=bookpurple,anchor=north] at (12,-3.0) {試行 4};
  % 白い丸の試行の番号（時刻 1 と 3 では丸が重なるため、下に書き出す）
  \node[font=\scriptsize,anchor=north east] at (-0.35,-11.75) {白い丸の試行};
  \def\rlbA{\textcolor{bookblue}{1}}\def\rlbB{\textcolor{bookorange}{2}}%
  \def\rlbC{\textcolor{bookgreen}{3}}\def\rlbD{\textcolor{bookpurple}{4}}%
  \foreach \t/\k in {1/{\rlbA\\\rlbB\\\rlbD},2/{\rlbD},3/{\rlbC\\\rlbD},4/{\rlbC},5/{\rlbC}}
    \node[font=\scriptsize,anchor=north,align=center,inner sep=1pt] at (\t,-11.75) {\k};
  \node[font=\scriptsize,anchor=north west] at (5.6,-11.75) {（計 8 個）};
  \node[font=\scriptsize,anchor=north west,align=left,draw=black!30,fill=white,inner sep=3pt]
    at (12.9,0.2) {\textbf{状態 9 で集めた $G_t$}\\\textcolor{bookblue}{試行 1}: $-0.81$\\
\textcolor{bookorange}{試行 2}: $-0.43$\\
\textcolor{bookgreen}{試行 3}: $+0.53$、$+0.59$、$+0.66$\\
\textcolor{bookpurple}{試行 4}: $+0.35$、$+0.39$、$+0.43$\\
8 個の平均: $+0.21$\\
正しい $V^\pi(9)=-0.155$
};
  \node[font=\scriptsize,anchor=north west,align=left] at (12.9,-6.6) {○ 状態 9 を通った時刻\\■ 終わり（ゴールか穴）};
\end{tikzpicture}

\medskip
\begin{tikzpicture}
\begin{axis}[height=0.3\linewidth,width=0.74\linewidth,xlabel={試行（エピソード）の数（対数目盛）},ylabel={$V^\pi(9)$ の見積もり},
  xmode=log,xmin=2,xmax=2000,ymin=-0.55,ymax=0.15,legend pos=outer north east,legend cell align=left]
\addplot[bookaccent!70,thin] table[col sep=comma,x=episode,y=alpha005]{code/ch20/out/fig_mc_spacetime.csv};
\addplot[bookblue,very thick] table[col sep=comma,x=episode,y=mean]{code/ch20/out/fig_mc_spacetime.csv};
\addplot[black,dashed,domain=2:2000,samples=2] {-0.155};
\legend{定数 $\alpha=0.05$,$1/n$ で平均,正しい値}
\end{axis}
\end{tikzpicture}
\caption{モンテカルロ法の時空図。上: 横軸が時刻、縦軸がマス。1 本の折れ線が 1 回の試行（ランダム方策）。
縦軸は状態の番号を順に並べただけで、縦の距離は迷路での距離ではない。
斜めの線は 1 歩の移動で、途中の行の状態は通っていない
（たとえば試行 1 の時刻 3→4 は、状態 11 から上の穴 7 への 1 歩で、状態 9 の帯を横切るが、状態 9 は通っていない）。
線が重ならないよう、試行ごとに少し上下にずらしたが、時刻 0〜3 の試行 1 と 2 は同じマスを通るので、線がほぼ重なる。
時刻 1 と 3 では白い丸も重なるので、丸のある試行の番号を、試行の色で時刻の目盛りの下に書いた（白い丸は全部で 8 個）。
図を見やすくするため、12 歩以内で終わり、状態 9 を通った試行を、先頭から順に 4 本選んで描いた。
状態 9 を通るたびに（白い丸）、そこから先の収益 $G_t$ を集め、平均して $V^\pi(9)$ の見積もりにする。
下: 試行を増やしたときの見積もり（上とは別の 2,000 本。選ばずに全部使った）。$1/n$ の平均は正しい値に近づき、定数 $\alpha$ の更新は上下し続ける。
横軸は対数目盛なので、右ほど 1 目盛りに多くの試行が詰まり、赤い線の上下が密に見える。数十本より先では、上下の幅はほぼ変わらない。
最初の 4 本では、定数 $\alpha$ の見積もりがたまたま $-0.153$ で、正しい値の破線とほぼ重なる。}
\label{fig:rlb-mc-spacetime}
\end{figure}

平均は、データが届くたびに少しずつ更新できます。
$n$ 個のデータ $x_1,\dots,x_n$ の平均を $\bar x_n$ と書きます。
\begin{align}
\bar x_n&=\frac1n\,(x_1+\cdots+x_{n-1}+x_n) \notag\\
&=\frac1n\,\bigl((n-1)\,\bar x_{n-1}+x_n\bigr) && \text{はじめの $n-1$ 個の和は $(n-1)\bar x_{n-1}$}\notag\\
&=\bar x_{n-1}+\frac1n\,(x_n-\bar x_{n-1}) && \text{展開して整理}
\label{eq:rlb-running-mean}
\end{align}
3 行目では、$\frac{n-1}{n}\,\bar x_{n-1}$ を $\bar x_{n-1}-\frac1n\,\bar x_{n-1}$ と分けてから、$\frac1n$ でくくりました。
新しい平均は、「古い平均」を「新しいデータとのずれ」の $1/n$ だけ動かしたものです。
数で確かめます。
2 と 4 の平均は 3 です。
そこへ 9 が届くと、新しい平均は $3+\frac13\,(9-3)=5$ で、計算し直した $(2+4+9)/3=5$ と同じです。

$1/n$ を小さな定数 $\alpha$ に置き換えると、モンテカルロ法の更新式になります。
定数にすると、古いデータの影響が少しずつ薄れるので、学習の途中で目標が変わっていく場合（このあとの TD 学習と Q 学習がそうです）にも追いかけられます。
その代わり、見積もりは最後まで少し上下し続けます。
\figref{fig:rlb-mc-spacetime}の下の図の赤い線が、$\alpha=0.05$ の場合です。
2,000 本の試行のあとでも、正しい値のまわりを上下しています。
\begin{equation}
V(s_t)\leftarrow V(s_t)+\alpha\,\bigl(G_t-V(s_t)\bigr)
\label{eq:rlb-mc}
\end{equation}
「$\leftarrow$」は、右辺の値で左辺を書き換えるという意味です。
$\alpha$ は、勾配降下法の学習率（\chapref{ch:ml}）と同じ役割の数で、ここでも学習率と呼びます。
これまでの章では $\eta$ と書きましたが、この章では、強化学習の慣例に合わせて $\alpha$ と書きます。
この式は、「見積もりを、目標 $G_t$ に少しだけ寄せる」と読めます。
この先の更新式は、目標が違うだけで、すべてこの形をしています。

モンテカルロ法には、弱点が 2 つあります。
エピソードが終わるまで $G_t$ が分からないので、途中では学習できません。
また、$G_t$ は多くの偶然が重なった結果なので、エピソードごとのばらつきが大きくなります。

\subsection*{TD 学習：1 歩先の見積もりを使う}

収益の再帰式 $G_t=r_t+\gamma G_{t+1}$ を思い出します。
$G_{t+1}$ はエピソードの最後まで分かりません。
そこで、その代わりに、\textbf{いま持っている見積もり} $V(s_{t+1})$ を使います。
すると、目標は $r_t+\gamma V(s_{t+1})$ となり、1 歩進むだけで計算できます。
この目標には、学習中の見積もり $V$ が入っています。
そのため、学習が進んで $V$ が変わると、目標も変わります。
モンテカルロ法の小節で、$1/n$ の代わりに定数 $\alpha$ を使う理由として挙げた「目標が変わっていく場合」が、これです。

\begin{definition}[TD 誤差と TD 学習]\label{def:rlb-td}
次の $\delta_t$（デルタ）を\term{TD 誤差}{TD error}と呼びます。
\begin{equation}
\delta_t\defeq r_t+\gamma V(s_{t+1})-V(s_t)
\label{eq:rlb-td-error}
\end{equation}
$s_{t+1}$ が終端状態のときは、$V(s_{t+1})=0$ とします。
TD 誤差を使って、1 歩ごとに
\begin{equation}
V(s_t)\leftarrow V(s_t)+\alpha\,\delta_t
\label{eq:rlb-td}
\end{equation}
と更新する方法を、\term{TD 学習}{temporal-difference learning}と呼びます。
\end{definition}

TD は temporal difference（時間的な差）の略です。
TD 誤差は、「1 歩進んでみたら、予想より良かったか、悪かったか」を表します。
1 歩前の予想が $V(s_t)$ で、1 歩進んでから立て直した予想が $r_t+\gamma V(s_{t+1})$ です。
$V$ が正しい $V^\pi$ に一致していれば、TD 誤差の期待値は 0 になります。
これは、ベルマン期待方程式の言い換えです。
ベルマン期待方程式は、「$r_t+\gamma V^\pi(s_{t+1})$ の平均は、$V^\pi(s_t)$ に等しい」という式でした。
両辺から $V^\pi(s_t)$ を引けば、「TD 誤差の平均は 0」になります。

\begin{intuition}{到着を待たずに、予想を直す}
家を出るときに、「会社まで 30 分」と予想したとします。
10 分歩いて駅に着くと、電車が遅れていて、「ここからあと 35 分」と分かりました。
この時点で、最初の予想は $10+35=45$ 分に直せます。
会社に着くまで待つ必要はありません。
モンテカルロ法は、着いてから実際の時間で直します。
TD 学習は、途中の予想を使って、1 歩ごとに直します。
式\eqref{eq:rlb-td-error}では、30 分が $V(s_t)$、10 分が $r_t$、35 分が $V(s_{t+1})$ に当たり、TD 誤差は $10+35-30=15$ 分です（$\gamma=1$ の場合）。
\end{intuition}

見積もりを使って見積もりを更新することを、\term{ブートストラップ}{bootstrapping}と呼びます。
\figref{fig:rlb-dp-mc-td}に、3 つの方法の違いをまとめます。

\begin{figure}[tbp]
\centering
\begin{tikzpicture}[
  st/.style={draw=bookblue,fill=figblue,circle,minimum size=3.6mm,inner sep=0pt},
  ac/.style={fill=black!70,circle,minimum size=1.8mm,inner sep=0pt},
  tm/.style={draw=black!60,fill=black!30,minimum size=3mm,inner sep=0pt},
  ln/.style={thick,black!65}]
  \begin{scope}
    \node[font=\small\bfseries] at (0,0.6) {動的計画法};
    \node[st] (s) at (0,0) {};
    \foreach \i/\x in {1/-0.9,2/0.9} {\node[ac] (a\i) at (\x,-0.8) {}; \draw[ln] (s)--(a\i);}
    \foreach \i/\x/\p in {1/-1.35/1,2/-0.45/1,3/0.45/2,4/1.35/2}
      {\node[st] (n\i) at (\x,-1.6) {}; \draw[ln] (a\p)--(n\i);}
    \node[lbl,align=center] at (0,-3.55) {1 歩先のすべての場合を\\確率で平均する};
  \end{scope}
  \begin{scope}[xshift=4.4cm]
    \node[font=\small\bfseries] at (0,0.6) {モンテカルロ法};
    \node[st] (m0) at (0,0) {};
    \node[ac] (m1) at (0,-0.5) {};
    \node[st] (m2) at (0,-1.0) {};
    \node[ac] (m3) at (0,-1.5) {};
    \node[st] (m4) at (0,-2.0) {};
    \node[ac] (m5) at (0,-2.5) {};
    \node[tm] (m6) at (0,-3.0) {};
    \foreach \a/\b in {m0/m1,m1/m2,m2/m3,m3/m4,m4/m5,m5/m6} \draw[ln] (\a)--(\b);
    \node[lbl,anchor=west] at (0.25,-3.0) {終端};
    \node[lbl,align=center] at (0,-3.55) {実際の 1 本の軌跡を\\最後までたどる};
  \end{scope}
  \begin{scope}[xshift=8.6cm]
    \node[font=\small\bfseries] at (0,0.6) {TD 学習};
    \node[st] (t0) at (0,0) {};
    \node[ac] (t1) at (0,-0.5) {};
    \node[st] (t2) at (0,-1.0) {};
    \draw[ln] (t0)--(t1)--(t2);
    \node[lbl,anchor=west] at (0.25,-1.0) {ここで $V$ を使う};
    \node[lbl,align=center] at (0,-3.55) {実際の 1 歩だけ進み、\\先は見積もりで代用する};
  \end{scope}
\end{tikzpicture}
\caption{価値の更新に使う情報の違い（模式図）。青い丸が状態、黒い点が行動、灰色の四角が終端状態。
動的計画法は、環境のモデルを使って 1 歩先のすべての場合を調べる（図では行動 2 つ、次の状態 2 つずつだけを描いた）。
モンテカルロ法と TD 学習は、実際に経験した 1 本の軌跡だけを使う。}
\label{fig:rlb-dp-mc-td}
\end{figure}

2 つの更新式を、1 エピソードだけ手で動かして比べます。
経験するのは、\figref{fig:rlb-trajectory}の軌跡（$8\to9\to10\to6\to2\to$ゴール）です。
表 $V$ はすべて 0 から始め、$\alpha=0.05$ とします。
結果が\tabref{tab:rlb-mc-td-hand}です。

\begin{table}[tbp]
\centering
\caption{同じ 1 本の軌跡（\figref{fig:rlb-trajectory}）で、0 から始めた表 $V$ を更新した結果（$\alpha=0.05$、$\gamma=0.9$）。
モンテカルロ法は通った状態のすべてを、TD 学習はゴールの手前の状態 2 だけを動かす。
更新後の値は、\figref{fig:rlb-mc-td-frames}と同じく小数第 4 位まで書いた（状態 8 は $0.05\times0.6561=0.032805$）。}
\label{tab:rlb-mc-td-hand}
\small
\begin{tabular}{ccllll}
\toprule
 & & \multicolumn{2}{c}{モンテカルロ法} & \multicolumn{2}{c}{TD 学習} \\
\cmidrule(lr){3-4}\cmidrule(lr){5-6}
時刻 $t$ & 状態 $s_t$ & 目標 $G_t$ & 更新後の $V(s_t)$ & 目標 $r_t+\gamma V(s_{t+1})$ & 更新後の $V(s_t)$ \\
\midrule
0 & 8 & 0.6561 & 0.0328 & 0 & 0.0000 \\
1 & 9 & 0.7290 & 0.0365 & 0 & 0.0000 \\
2 & 10 & 0.8100 & 0.0405 & 0 & 0.0000 \\
3 & 6 & 0.9000 & 0.0450 & 0 & 0.0000 \\
4 & 2 & 1.0000 & 0.0500 & 1 & 0.0500 \\
\bottomrule
\end{tabular}
\end{table}

モンテカルロ法の目標は、エピソードが終わってから計算した収益 $G_t$ です。
たとえば状態 8 は、$0.05\times0.6561\approx0.033$ になります。
通った 5 つの状態のすべてで、ゴールの報酬の影響を受けて、値が 0 から動きます。
TD 学習では、はじめの 4 歩は、報酬も次の状態の見積もりも 0 なので、TD 誤差が 0 で、表は変わりません。
最後の 1 歩だけ、TD 誤差が $1+0.9\times0-0=1$ になり、状態 2 が $0.05$ に変わります。
同じ軌跡をもう一度経験すると、状態 6 の TD 誤差が $0.9\times0.05=0.045$ になり、値は $0.05\times0.045=0.00225$ に変わります。
TD 学習では、ゴールの報酬の影響は、1 エピソードに 1 マスずつしか手前へ伝わりません（「伝わる」は、そのマスの値が 0 から動くことです）。

\figref{fig:rlb-mc-td-frames}は、同じ軌跡を 3 回くり返し経験したときの表を、1 回ごとに並べたコマ送りの図です（\file{ch20/fig_dp_extra.py}）。
赤い枠が、その回に値が変わったマスです。
左のモンテカルロ法では、1 回目から 5 マスすべてに枠があり、回を重ねるごとに全部のマスが同じように増えます。
右の TD 学習では、枠のあるマスが、1 回目は状態 2、2 回目は状態 6 と 2、3 回目は状態 10 と 6 と 2 と、1 回に 1 マスずつ左へ広がります。
状態 8 と 9 は、3 回目でもまだ 0 です。
どちらが速く正しい値に近づくのかは、次の小節で実際に測ります。

\begin{figure}[tbp]
\centering
\begin{tikzpicture}[x=1.05cm,y=0.62cm]
  \foreach \j/\s in {0/8,1/9,2/10,3/6,4/2,6/8,7/9,8/10,9/6,10/2}
    \node[font=\scriptsize] at (\j+0.5,0.35) {状態 \s};
  \node[font=\small\bfseries] at (2.5,1.05) {モンテカルロ法};
  \node[font=\small\bfseries] at (8.5,1.05) {TD 学習};
  \foreach \k in {1,2,3}
    \node[font=\scriptsize,anchor=east] at (-0.1,-\k+0.5) {\k 回目};
  \rlbft{0}{0}{0.0328}{bookgreen!13}{1}
\rlbft{0}{1}{0.0365}{bookgreen!15}{1}
\rlbft{0}{2}{0.0405}{bookgreen!16}{1}
\rlbft{0}{3}{0.0450}{bookgreen!18}{1}
\rlbft{0}{4}{0.0500}{bookgreen!20}{1}
\rlbft{0}{6}{0.0000}{bookgreen!0}{0}
\rlbft{0}{7}{0.0000}{bookgreen!0}{0}
\rlbft{0}{8}{0.0000}{bookgreen!0}{0}
\rlbft{0}{9}{0.0000}{bookgreen!0}{0}
\rlbft{0}{10}{0.0500}{bookgreen!20}{1}
\rlbft{1}{0}{0.0640}{bookgreen!26}{1}
\rlbft{1}{1}{0.0711}{bookgreen!28}{1}
\rlbft{1}{2}{0.0790}{bookgreen!32}{1}
\rlbft{1}{3}{0.0878}{bookgreen!35}{1}
\rlbft{1}{4}{0.0975}{bookgreen!39}{1}
\rlbft{1}{6}{0.0000}{bookgreen!0}{0}
\rlbft{1}{7}{0.0000}{bookgreen!0}{0}
\rlbft{1}{8}{0.0000}{bookgreen!0}{0}
\rlbft{1}{9}{0.0023}{bookgreen!1}{1}
\rlbft{1}{10}{0.0975}{bookgreen!39}{1}
\rlbft{2}{0}{0.0936}{bookgreen!37}{1}
\rlbft{2}{1}{0.1040}{bookgreen!42}{1}
\rlbft{2}{2}{0.1155}{bookgreen!46}{1}
\rlbft{2}{3}{0.1284}{bookgreen!51}{1}
\rlbft{2}{4}{0.1426}{bookgreen!57}{1}
\rlbft{2}{6}{0.0000}{bookgreen!0}{0}
\rlbft{2}{7}{0.0000}{bookgreen!0}{0}
\rlbft{2}{8}{0.0001}{bookgreen!0}{1}
\rlbft{2}{9}{0.0065}{bookgreen!3}{1}
\rlbft{2}{10}{0.1426}{bookgreen!57}{1}

  \draw[black!45,step=1] (0,-3) grid (5,0);
  \draw[black!45,step=1] (6,-3) grid (11,0);
\end{tikzpicture}
\caption{同じ軌跡（$8\to9\to10\to6\to2\to$ゴール）を 3 回くり返し経験したときの表 $V$（$\alpha=0.05$、$\gamma=0.9$）。1 行が 1 回分のコマ。
赤い枠は、その回に値が変わったマス。緑が濃いほど値が大きい。
モンテカルロ法は毎回すべてのマスを動かし、TD 学習は 1 回に 1 マスずつゴールから遠くへ広げる。1 回目の行は\tabref{tab:rlb-mc-td-hand}と同じ値。}
\label{fig:rlb-mc-td-frames}
\end{figure}

\subsection*{実装して比べる}
\skimread{押さえる結論は 2 つです。この迷路では、学習率が小さめのときは、はじめはモンテカルロ法が速く、最後は TD 学習のほうが小さい誤差まで下がります。モンテカルロ法の目標はばらつきが大きく、TD 学習の目標は、ばらつきが小さい代わりに偏りがあります。}

\snippet{モンテカルロ法と TD 学習の更新。違いは「何に寄せるか」だけ}{code:rlb-mc-td}{ch20/mc_td.py}{mc-td}

どちらの関数も、1 エピソード分の軌跡 \code{traj} を受け取って、表 \code{V} を書き換えます。
\code{mc_update} の 5 行目が式\eqref{eq:rlb-mc}で、目標は実際の収益 \code{G[t]} です。
\code{td_update} の 13 行目が TD 誤差（式\eqref{eq:rlb-td-error}）、14 行目が式\eqref{eq:rlb-td}です。
12 行目では、最後の 1 歩のときに、次の状態（終端状態）の価値を 0 にしています。
どちらも、遷移確率の表 \code{P} を一度も使っていません。
なお、\code{td_update} もエピソードのあとに呼ぶのは、2 つの方法に同じ軌跡を与えて比べるためです。
\code{for} 文の中身はその 1 歩の情報しか使わないので、歩きながら実行しても同じ結果になります。

ランダム方策で 300 エピソードを動かしました。
\figref{fig:rlb-mc-td}は、\figref{fig:rlb-v-random}の正しい $V^\pi$ との誤差です（100 回の実行の平均）。

\begin{figure}[tbp]
\centering
\begin{tikzpicture}
\begin{axis}[at={(0\linewidth,0)},height=0.3\linewidth,width=0.355\linewidth,
  xlabel={エピソード数},ylabel={正しい $V^\pi$ との誤差},
  title={\small $\alpha=0.02$},xmin=0,xmax=300,ymin=0,ymax=0.3,xtick={0,100,200,300},legend style={font=\scriptsize,at={(1.45,1.22)},anchor=south,draw=none,/tikz/every even column/.append style={column sep=5mm}},legend columns=2,legend cell align=left]
\addplot[bookaccent,thick] table[col sep=comma,x=episode,y=mc002]{code/ch20/out/mc_td.csv};
\addplot[bookblue,thick,dashed] table[col sep=comma,x=episode,y=td002]{code/ch20/out/mc_td.csv};
\legend{モンテカルロ法,TD 学習}
\end{axis}
\begin{axis}[at={(0.32\linewidth,0)},height=0.3\linewidth,width=0.355\linewidth,
  xlabel={エピソード数},
  title={\small $\alpha=0.05$},xmin=0,xmax=300,ymin=0,ymax=0.3,xtick={0,100,200,300},yticklabels={}]
\addplot[bookaccent,thick] table[col sep=comma,x=episode,y=mc005]{code/ch20/out/mc_td.csv};
\addplot[bookblue,thick,dashed] table[col sep=comma,x=episode,y=td005]{code/ch20/out/mc_td.csv};
\end{axis}
\begin{axis}[at={(0.64\linewidth,0)},height=0.3\linewidth,width=0.355\linewidth,
  xlabel={エピソード数},
  title={\small $\alpha=0.1$},xmin=0,xmax=300,ymin=0,ymax=0.3,xtick={0,100,200,300},yticklabels={}]
\addplot[bookaccent,thick] table[col sep=comma,x=episode,y=mc01]{code/ch20/out/mc_td.csv};
\addplot[bookblue,thick,dashed] table[col sep=comma,x=episode,y=td01]{code/ch20/out/mc_td.csv};
\end{axis}

\end{tikzpicture}
\caption{モンテカルロ法（赤の実線）と TD 学習（青の破線）の誤差（9 状態の誤差を 2 乗して平均し、平方根を取った値。100 回の実行の平均）。
学習率 $\alpha$ だけを 3 通りに変え、それ以外（ランダム方策、$\gamma=0.9$、エピソード）は同じ。
$\alpha$ を大きくすると、どちらも速く下がるが、下げ止まる高さが上がる。
$\alpha=0.02$ と $0.05$ では、はじめはモンテカルロ法が速く、TD 学習があとで追い越す。$\alpha=0.1$ では、10 エピソードより前から TD 学習のほうが小さい。}
\label{fig:rlb-mc-td}
\end{figure}

学習率 $\alpha=0.05$ の結果（\figref{fig:rlb-mc-td}の中央）を見ます。
はじめの 10 エピソードでは、モンテカルロ法の誤差（0.169）のほうが、TD 学習（0.208）より小さくなりました。
モンテカルロ法は、1 エピソードで、通ったすべての状態の値に最後の結果を反映させるからです（\figref{fig:rlb-mc-td-frames}）。
TD 学習では、ゴールの報酬は 1 回に 1 マスずつしか伝わりません（\tabref{tab:rlb-mc-td-hand}で見たとおりです）。
しかし、モンテカルロ法の誤差は 0.13 付近で下げ止まります。
収益のばらつきが大きいうえに、$\alpha$ を定数にしているので、新しいエピソードのたびに表の値が上下し続けるためです（\figref{fig:rlb-mc-spacetime}の下の図の赤い線）。
$1/n$ で平均すれば、この上下は小さくなっていきます（同じ図の青い線）。
TD 学習は、300 エピソードで 0.052 まで下がりました。

学習率を変えると、結果はどう変わるでしょうか。
\figref{fig:rlb-mc-td}の 3 枚を、左から順に比べてください（実行結果 \file{ch20/out/mc_td.txt}）。
$\alpha$ を大きくすると、どちらの方法も速く下がりますが、下げ止まる高さが上がります。
300 エピソード目の誤差は、モンテカルロ法が 0.077、0.128、0.187、TD 学習が 0.032、0.052、0.076 です。
1 回の更新で大きく動かすほど、1 本のエピソードの偶然の影響を強く受けるからです。
TD 学習がモンテカルロ法を追い越すのは、$\alpha=0.02$ では 138 エピソード目、$0.05$ では 33 エピソード目、$0.1$ では 8 エピソード目です。
どの $\alpha$ でも、最後は TD 学習のほうが誤差が小さくなりました。
ただし、これはこの問題での結果です。いつでも TD 学習が勝つわけではありません。

2 つの方法の違いは、目標の性質にあります。
モンテカルロ法の目標 $G_t$ は、平均すれば正しい値になりますが（偏りがない）、偶然が何歩分も重なるので、毎回大きく違います（ばらつきが大きい）。
TD 学習の目標は、偶然が 1 歩分しか入らないので、ばらつきは小さくなります。
その代わり、まだ正しくない見積もり $V(s_{t+1})$ を含むので、平均しても正しい値からずれています（偏りがある）。
\tabref{tab:rlb-dp-mc-td}に、3 つの方法の性質をまとめます。

\begin{table}[tbp]
\centering
\caption{価値を求める 3 つの方法。「偏り」は、目標の平均が正しい値からずれること。動的計画法の目標も、途中の見積もり $V$ を含む。}
\label{tab:rlb-dp-mc-td}
\small
\begin{tabular}{llll}
\toprule
 & 動的計画法 & モンテカルロ法 & TD 学習 \\
\midrule
環境のモデル & 必要 & 不要 & 不要 \\
更新の目標 & 全部の場合の平均 & 実際の収益 $G_t$ & $r_t+\gamma V(s_{t+1})$ \\
いつ更新できるか & 経験なしで（表 $P$ だけで） & エピソードのあと & 1 歩ごと \\
目標のばらつき & なし & 大きい & 小さい \\
目標の偏り & 見積もりの分だけ & なし & 見積もりの分だけ \\
\bottomrule
\end{tabular}
\end{table}

この節で分かったことをまとめます。
環境のモデルがなくても、経験した収益を平均すれば価値を推定できます（モンテカルロ法）。
TD 学習は、1 歩先の見積もりを目標に使い、エピソードの終わりを待たずに学習します。
次章でも、TD 誤差を使います。

\section{Q 学習と、探索と活用}\label{sec:rlb-q}

TD 学習でできたのは、決まった方策の価値を測ることでした。
この節の問いは、「環境のモデルを知らないまま、最適な方策そのものを学ぶにはどうするか」です。

\subsection*{Q 学習}

方策を良くするには、各状態で、価値が最大の行動を知る必要があります。
$V$ から行動の価値を計算する式\eqref{eq:rlb-q-from-v}は、遷移確率 $P$ を使います。
環境のモデルがないときは、この式が使えません。
そこで、はじめから行動価値 $Q(s,a)$ の表を学習します。

目標にするのは、$Q^\ast$ のベルマン最適方程式\eqref{eq:rlb-bellman-opt-q}です。
TD 学習と同じように、次の状態についての平均を、実際に経験した 1 回分で代用します。
\begin{equation}
Q(s_t,a_t)\leftarrow Q(s_t,a_t)+\alpha\,\bigl[\,r_t+\gamma\max_{a'}Q(s_{t+1},a')-Q(s_t,a_t)\,\bigr]
\label{eq:rlb-q-learning}
\end{equation}
この方法が\term{Q 学習}{Q-learning}です\cite{watkins1992q}。
状態と行動の組ごとに行動価値の数を並べた表を持ち、1 歩ごとに 1 か所を書き換えます。

数で追ってみます。表はすべて 0 で、$\alpha=0.1$ とします。
エージェントが、状態 2 で「右」を選んで、初めてゴールに入ったとします。
報酬は 1 で、終端状態の $Q$ は 0 なので、目標は $1+0.9\times0=1$ です。
$Q(2,\text{右})$ は、$0+0.1\times(1-0)=0.1$ になります。
次のエピソードで、状態 1 から「右」を選んで、状態 2 に進んだとします。
報酬は 0 ですが、状態 2 にはもう $0.1$ が入っているので、目標は $0.9\times0.1=0.09$ です。
$Q(1,\text{右})$ は、$0.1\times0.09=0.009$ になります。
価値反復と同じように、ゴールの報酬が 1 マスずつ手前へ伝わります。
違うのは、遷移確率の表を使わず、実際に歩いた 1 歩だけで書き換えることです。

すべての組を十分に何度も経験し、学習率を適切に小さくしていけば、表は $Q^\ast$ に収束することが証明されています\cite{watkins1992q}。

ここで、新しい問題が生まれます。
「すべての組を経験する」ためには、エージェントはどう行動を選べばよいのでしょうか。

\subsection*{探索と活用：多腕バンディット}

この問題だけを取り出して考えるために、いちばん簡単な設定を使います。
スロットマシンが 5 台あり、当たる確率は台ごとに違います（0.2、0.8、0.5、0.3、0.6）。
当たれば報酬 1、外れれば 0 です。
プレイヤーは、確率を知りません。
状態は 1 つだけで、1 回引いたら終わりです（1 回引くことが、1 エピソードに当たります）。
この問題を\term{多腕バンディット}{multi-armed bandit}と呼びます。

各台の価値 $Q(a)$ は、これまでに引いた結果の平均で見積もれます。
見積もりが最大の台を引き続けるのが、いままでの知識の\term{活用}{exploitation}です。
しかし、活用だけでは困ったことが起きます。
最初にたまたま当たった台の見積もりが高くなり、ほかの台を二度と試さなくなるのです。
もっと良い台を見つけるには、見積もりが低い台もあえて引く\textbf{探索}が必要です。
いちばん簡単な方法が、\term{$\varepsilon$-貪欲法}{$\varepsilon$-greedy}です。
確率 $\varepsilon$（イプシロン）でランダムに選び、残りの確率 $1-\varepsilon$ で見積もりが最大のものを選びます。

\snippet{$\varepsilon$-貪欲法}{code:rlb-bandit}{ch20/bandit.py}{eps-greedy}

3〜4 行目が探索、5〜6 行目が活用です。
見積もりが最大の台が複数あるときは、その番号をすべて取り出し（5 行目の \code{np.flatnonzero}）、その中からランダムに選びます。
この関数で台を選び、結果が出るたびに、その台の見積もりを式\eqref{eq:rlb-running-mean}で更新します（当たる確率は変わらないので、$1/n$ のまま使います）。
これを 500 回繰り返しました（同じファイルの \code{run_bandit}）。

\begin{figure}[tbp]
\centering
\begin{tikzpicture}
\begin{axis}[height=0.33\linewidth,width=0.62\linewidth,xlabel={引いた回数},ylabel={最良の台を選んだ割合},
  xmin=0,xmax=500,ymin=0,ymax=1,legend pos=outer north east,legend cell align=left]
\addplot[bookgray,dashed,thick] table[col sep=comma,x=step,y=b0]{code/ch20/out/bandit.csv};
\addplot[bookblue] table[col sep=comma,x=step,y=b01]{code/ch20/out/bandit.csv};
\addplot[bookaccent] table[col sep=comma,x=step,y=b05]{code/ch20/out/bandit.csv};
\legend{$\varepsilon=0$（活用だけ）,$\varepsilon=0.1$,$\varepsilon=0.5$}
\end{axis}
\end{tikzpicture}
\caption{多腕バンディットでの $\varepsilon$-貪欲法（2,000 回の実行の平均）。
探索しないと最良の台を見つけられず、探索しすぎると見つけた台を活かせない。}
\label{fig:rlb-bandit}
\end{figure}

\figref{fig:rlb-bandit}が結果です。
$\varepsilon=0$ では、最良の台（当たる確率 0.8）を選ぶ割合が 31\% で止まります。
最初に当たりが出た台の見積もりだけが正になり、あとはその台を引き続けるからです。
外れた台の見積もりは 0 のままで、ほかの台と同点です。
そのため、最初の当たりが出るまでは、毎回 5 台から等しい確率で選びます。
すると、最初の当たりがどの台で出るかは、各台の当たる確率に比例します。
それが最良の台である確率は $0.8/(0.2+0.8+0.5+0.3+0.6)=1/3$ です。
実測の 31\% は、これに近い値です。
$\varepsilon=0.1$ では、500 回目に 88\% まで上がります。
$\varepsilon=0.5$ では、早く見つけますが、半分はでたらめに引くので、60\% で頭打ちです（でたらめに引いた半分のうち、5 回に 1 回は最良の台に当たります）。
500 回の平均報酬は、順に 0.563、0.729、0.633 でした。
探索は、少なすぎても多すぎても損をします。
これを\jterm{探索と活用のトレードオフ}と呼びます。

\subsection*{Q 学習を実装する}

\snippet{表形式の Q 学習}{code:rlb-q-learning}{ch20/q_learning.py}{q-learning}

7 行目で、バンディットと同じ $\varepsilon$-貪欲法で行動を選びます。
\code{Q[s]} は、状態 $s$ での 4 つの行動の価値です。
8 行目で 1 歩進み、9〜10 行目が式\eqref{eq:rlb-q-learning}です。
終端状態の行は一度も書き換えられず 0 のままなので、最後の 1 歩の目標は報酬だけになります。
12〜13 行目は、学習の途中経過を記録するための仕掛けです。

$\alpha=0.1$、$\varepsilon=0.1$ で 2,000 エピソード学習した結果が、\figref{fig:rlb-q-table}です。

\begin{figure}[tbp]
\centering
\begin{tikzpicture}[x=1.5cm,y=1.35cm]
  \renewcommand{\rlbfont}{\tiny}
  \rlbq{0}{0}{0}{0.66}{bookgreen!46}{0}
\rlbq{0}{0}{1}{0.75}{bookgreen!52}{1}
\rlbq{0}{0}{2}{0.60}{bookgreen!42}{0}
\rlbq{0}{0}{3}{0.64}{bookgreen!45}{0}
\rlbq{1}{0}{0}{0.72}{bookgreen!50}{0}
\rlbq{1}{0}{1}{0.85}{bookgreen!60}{1}
\rlbq{1}{0}{2}{0.74}{bookgreen!52}{0}
\rlbq{1}{0}{3}{0.66}{bookgreen!46}{0}
\rlbq{2}{0}{0}{0.81}{bookgreen!57}{0}
\rlbq{2}{0}{1}{0.99}{bookgreen!69}{1}
\rlbq{2}{0}{2}{0.51}{bookgreen!35}{0}
\rlbq{2}{0}{3}{0.68}{bookgreen!48}{0}
\rlbgoal{3}{0}
\rlbq{0}{1}{0}{0.66}{bookgreen!46}{1}
\rlbq{0}{1}{1}{0.57}{bookgreen!40}{0}
\rlbq{0}{1}{2}{0.52}{bookgreen!36}{0}
\rlbq{0}{1}{3}{0.56}{bookgreen!39}{0}
\rlbwall{1}{1}
\rlbq{2}{1}{0}{0.42}{bookgreen!30}{1}
\rlbq{2}{1}{1}{$-$0.61}{bookaccent!43}{0}
\rlbq{2}{1}{2}{0.06}{bookgreen!4}{0}
\rlbq{2}{1}{3}{0.36}{bookgreen!25}{0}
\rlbpit{3}{1}
\rlbq{0}{2}{0}{0.59}{bookgreen!41}{1}
\rlbq{0}{2}{1}{0.43}{bookgreen!30}{0}
\rlbq{0}{2}{2}{0.46}{bookgreen!32}{0}
\rlbq{0}{2}{3}{0.49}{bookgreen!34}{0}
\rlbq{1}{2}{0}{0.34}{bookgreen!24}{0}
\rlbq{1}{2}{1}{0.15}{bookgreen!11}{0}
\rlbq{1}{2}{2}{0.25}{bookgreen!17}{0}
\rlbq{1}{2}{3}{0.49}{bookgreen!34}{1}
\rlbq{2}{2}{0}{0.27}{bookgreen!19}{1}
\rlbq{2}{2}{1}{0.00}{bookgreen!0}{0}
\rlbq{2}{2}{2}{0.07}{bookgreen!5}{0}
\rlbq{2}{2}{3}{0.03}{bookgreen!2}{0}
\rlbq{3}{2}{0}{$-$0.05}{bookaccent!3}{0}
\rlbq{3}{2}{1}{0.00}{bookgreen!0}{1}
\rlbq{3}{2}{2}{0.00}{bookgreen!0}{0}
\rlbq{3}{2}{3}{$-$0.04}{bookaccent!3}{0}

  \rlbframe
  \rlbnumstag
  \node[font=\small] at (2,0.2) {Q 学習で学んだ $Q$};
  \begin{scope}[shift={(4.25,0)}]
    \rlbq{0}{0}{0}{0.65}{bookgreen!46}{0}
\rlbq{0}{0}{1}{0.72}{bookgreen!50}{1}
\rlbq{0}{0}{2}{0.59}{bookgreen!41}{0}
\rlbq{0}{0}{3}{0.64}{bookgreen!45}{0}
\rlbq{1}{0}{0}{0.74}{bookgreen!52}{0}
\rlbq{1}{0}{1}{0.83}{bookgreen!58}{1}
\rlbq{1}{0}{2}{0.74}{bookgreen!52}{0}
\rlbq{1}{0}{3}{0.66}{bookgreen!47}{0}
\rlbq{2}{0}{0}{0.85}{bookgreen!60}{0}
\rlbq{2}{0}{1}{0.94}{bookgreen!66}{1}
\rlbq{2}{0}{2}{0.63}{bookgreen!44}{0}
\rlbq{2}{0}{3}{0.74}{bookgreen!52}{0}
\rlbgoal{3}{0}
\rlbq{0}{1}{0}{0.63}{bookgreen!44}{1}
\rlbq{0}{1}{1}{0.57}{bookgreen!40}{0}
\rlbq{0}{1}{2}{0.51}{bookgreen!35}{0}
\rlbq{0}{1}{3}{0.57}{bookgreen!40}{0}
\rlbwall{1}{1}
\rlbq{2}{1}{0}{0.64}{bookgreen!44}{1}
\rlbq{2}{1}{1}{$-$0.67}{bookaccent!47}{0}
\rlbq{2}{1}{2}{0.34}{bookgreen!24}{0}
\rlbq{2}{1}{3}{0.59}{bookgreen!41}{0}
\rlbpit{3}{1}
\rlbq{0}{2}{0}{0.55}{bookgreen!38}{1}
\rlbq{0}{2}{1}{0.45}{bookgreen!32}{0}
\rlbq{0}{2}{2}{0.48}{bookgreen!34}{0}
\rlbq{0}{2}{3}{0.50}{bookgreen!35}{0}
\rlbq{1}{2}{0}{0.44}{bookgreen!31}{0}
\rlbq{1}{2}{1}{0.47}{bookgreen!33}{0}
\rlbq{1}{2}{2}{0.44}{bookgreen!31}{0}
\rlbq{1}{2}{3}{0.48}{bookgreen!34}{1}
\rlbq{2}{2}{0}{0.53}{bookgreen!37}{1}
\rlbq{2}{2}{1}{0.33}{bookgreen!23}{0}
\rlbq{2}{2}{2}{0.45}{bookgreen!32}{0}
\rlbq{2}{2}{3}{0.45}{bookgreen!31}{0}
\rlbq{3}{2}{0}{$-$0.72}{bookaccent!51}{0}
\rlbq{3}{2}{1}{0.15}{bookgreen!10}{0}
\rlbq{3}{2}{2}{0.30}{bookgreen!21}{0}
\rlbq{3}{2}{3}{0.31}{bookgreen!22}{1}

    \rlbframe
    \rlbnumstag
    \node[font=\small] at (2,0.2) {正しい $Q^\ast$（価値反復）};
  \end{scope}
\end{tikzpicture}
\caption{学習した行動価値（左）と、価値反復で求めた正解（右）。各マスの 4 つの三角形が、上・右・下・左の行動の価値（負の値は赤）。太字は最大の行動。
よく通る上の道では正解に近いが、ほとんど通らない右下では学習が進んでいない。}
\label{fig:rlb-q-table}
\end{figure}

学習した $Q$ から作った貪欲方策は、9 状態のうち 8 状態で最適方策と同じ行動を選びます（違うのは、右下の状態 11 だけです）。
スタートから貪欲方策で動いたときの価値は $0.545$ で、最適方策と同じです。
エージェントは、遷移確率を一度も見ずに、最適方策と同じ進み方を見つけました。

一方で、$Q$ の値は、場所によって正確さが大きく違います。
上の道（状態 8、4、0、1、2）では、$Q^\ast$ に近い値です。
状態 11 では、値がほとんど 0 のままです。状態 10 でも、最大の「上」が 0.27 で、正解（0.53）の半分ほどです。どちらも $Q^\ast$ から大きく外れています（差は最大で 0.675）。
エージェントが上の道ばかり通り、右下にはめったに行かないからです。
\secref{sec:rlb-what}で述べた「自分の行動がデータを変える」が、ここに現れています。

\figref{fig:rlb-q-curve}は、$\varepsilon$ を変えたときの学習の進み方です。
$\varepsilon=0$ では、最初に見つけた道から離れられず、平均 0.24 で止まります。
$\varepsilon=0.1$ と $0.3$ では、最適な価値 0.545 に近づきます。
ただし、この図は、探索をやめて貪欲に動いたときの成績です。
探索しながら得た成績を測った\figref{fig:rlb-bandit}とは違い、学習中にでたらめな行動で損をした分は、図に現れていません。

\begin{figure}[tbp]
\centering
\begin{tikzpicture}
\begin{axis}[height=0.33\linewidth,xlabel={エピソード数},ylabel={貪欲方策の価値 $V(8)$},
  xmin=0,xmax=2000,ymin=0.15,ymax=0.6,legend style={at={(0.98,0.5)},anchor=east},legend columns=2]
\addplot[bookgray,dashed,thick] table[col sep=comma,x=episode,y=eps0]{code/ch20/out/q_learning.csv};
\addplot[bookblue] table[col sep=comma,x=episode,y=eps01]{code/ch20/out/q_learning.csv};
\addplot[bookaccent] table[col sep=comma,x=episode,y=eps03]{code/ch20/out/q_learning.csv};
\addplot[black,dashed,domain=0:2000,samples=2] {0.5452};
\legend{$\varepsilon=0$,$\varepsilon=0.1$,$\varepsilon=0.3$,最適 $V^\ast(8)$}
\end{axis}
\end{tikzpicture}
\caption{Q 学習の学習曲線（20 回の実行の平均）。縦軸は、その時点の $Q$ から作った貪欲方策（探索なし）を、遷移確率の表を使って正確に評価した値。}
\label{fig:rlb-q-curve}
\end{figure}

\begin{quiz}
$\varepsilon=0$ の Q 学習は、なぜ最初に見つけた道から離れられないのでしょうか。
$Q$ の表が 0 から始まることと、報酬が 0 以上の道を一度見つけたあとのことを考えてください。
\end{quiz}

ゴールに一度着くと、その道の上の $Q$ は正の値になります。
試していない行動の $Q$ は 0 のままなので、貪欲に選ぶかぎり、正の値を持つ同じ道だけを選び続けます。
もっと良い道があっても、試さないので気づけません。

なお、Q 学習のエージェントは、歩くときにはときどきでたらめに動きます（$\varepsilon$-貪欲）。
しかし、表に書き込む目標は、$\max$ を使った「このあと最善の行動を取り続けたら」という値です。
このように、データを集める方策と学習の対象の方策が違う方法を、\term{オフポリシー}{off-policy}の方法と呼びます。
次章では、これとは逆に、データを集めた方策そのものを良くしていく方法が中心になります。このような方法を、オンポリシー（on-policy）の方法と呼びます（\chapref{ch:policy-gradient}）。

ここまでに出てきた方法を、\tabref{tab:rlb-methods}に整理します。
行（横の並び）は環境のモデルを使うかどうか、列は何を求めるかの違いです。
右下の Q 学習が、この章の到達点です。

\begin{table}[tbp]
\centering
\caption{この章の方法の整理。上の行は遷移確率の表を使い（本章では、モデルが与えられた場合の動的計画法）、下の行は経験だけを使う。
右の列の方法は、どれも「価値が最大の行動を選ぶ」操作を含む。}
\label{tab:rlb-methods}
\small
\begin{tabular}{lll}
\toprule
 & 決めた方策の価値を測る & 最適な方策を求める \\
\midrule
環境のモデルを使う（モデルベース） & 反復方策評価 & 方策反復、価値反復 \\
経験だけから学ぶ（モデルフリー） & モンテカルロ法、TD 学習 & Q 学習 \\
\bottomrule
\end{tabular}
\end{table}

この節で分かったことをまとめます。
Q 学習は、行動価値の表を 1 歩ごとに書き換えて、環境のモデルなしで最適方策を学びます。
そのためには探索が必要で、探索と活用のつり合いが学習の結果を左右します。

\section{表から関数へ、そして LLM へ}\label{sec:rlb-fa}

格子世界の $Q$ の表は、9 状態 $\times$ 4 行動で、たった 36 個の数でした。
この節の問いは、「状態が多すぎて表を作れないときはどうするか。LLM には、どの方法が向いているか」です。

\subsection*{表の限界と関数近似}

表を使う方法には、2 つの限界があります。
1 つ目は、大きさです。
囲碁の局面は約 $10^{170}$ 通りあります。
ミニGPT の語彙は 4,096 種類です。
状態を「ここまでのトークン列」とすると、長さ 10 の列だけで $4096^{10}\approx10^{36}$ 通りになります。
これだけの数を並べた表は、計算機のメモリに置けません。
2 つ目は、\textbf{経験していない状態については何も分からない}ことです。
\figref{fig:rlb-q-table}の右下がそうでした。
表の各マスは独立なので、隣のマスで学んだことが役に立ちません。

そこで、表の代わりに、パラメータを持つ関数で価値を表します。
これを\term{関数近似}{function approximation}と呼びます。
関数にニューラルネットワークを使えば、似た状態には似た値が出るので、経験していない状態にも見当がつきます。

\subsection*{表は、いちばん単純なネットワーク}
\skimread{押さえる結論は 2 つです。$Q$ の表は、one-hot ベクトルを入れる線形層と同じものです。その線形層を深いネットワークに取り替えて Q 学習をするのが、DQN です。}

実は、表は「いちばん単純なニューラルネットワーク」と見ることができます（\figref{fig:rlb-table-to-net}）。
状態 $s$ を、$s$ 番目だけが 1 の one-hot ベクトルで表します（\chapref{ch:ml}の one-hot 表現と同じです）。
これを、バイアスなしの線形層に入れます。
線形層は、\chapref{ch:nn}の全結合層（\code{nn.Linear}）のことです。
線形層の重み行列を $\mW$ と書きます。形は 4 行 $\times$ 12 列（行動の数 $\times$ 状態の数）です。
入力は $s$ 番目だけが 1 なので、$\mW$ を掛けると、$s$ 列目だけが残ります。
出力は、$\mW$ の $s$ 列目、つまり「表の $s$ 行目」そのものです。
$\mW$ は、表を転置したもの（行と列を入れ替えたもの）になっています。
\figref{fig:rlb-table-to-net}では、$s=8$ の場合を、左右の橙のマスで示しました。

\begin{figure}[tbp]
\centering
\begin{tikzpicture}
  % 左: Q の表（12 行 × 4 列。壁と終端状態の行は使わない）
  \node[lbl] at (0.8,0.75) {表 $Q$（12 行 $\times$ 4 列）};
  \foreach \c/\lab in {0/上,1/右,2/下,3/左}
    \node[font=\tiny] at (0.2+0.4*\c,0.22) {\lab};
  \foreach \r in {0,...,11} {
    \node[font=\tiny,anchor=east] at (-0.03,-0.25*\r-0.125) {\r};
    \foreach \c in {0,...,3}
      \draw[black!45,fill=white] (0.4*\c,-0.25*\r) rectangle ++(0.4,-0.25);}
  \foreach \r in {3,5,7} \foreach \c in {0,...,3}
    \draw[black!45,fill=black!12] (0.4*\c,-0.25*\r) rectangle ++(0.4,-0.25);
  \foreach \c in {0,...,3}
    \draw[black!60,fill=figorange] (0.4*\c,-2.0) rectangle ++(0.4,-0.25);
  \node[font=\tiny,anchor=west] at (1.65,-2.125) {← 状態 8 の行};
  % 右: 重み行列 W（4 行 × 12 列）と one-hot ベクトル x
  \begin{scope}[xshift=3.3cm]
    \node[lbl] at (1.8,0.75) {線形層: $\mW\vx$};
    \foreach \cc in {0,...,11} {
      \node[font=\tiny] at (0.3*\cc+0.15,-0.7) {\cc};
      \foreach \rr in {0,...,3}
        \draw[black!45,fill=white] (0.3*\cc,-0.9-0.3*\rr) rectangle ++(0.3,-0.3);}
    \foreach \rr in {0,...,3}
      \draw[black!60,fill=figorange] (2.4,-0.9-0.3*\rr) rectangle ++(0.3,-0.3);
    \foreach \rr/\lab in {0/上,1/右,2/下,3/左}
      \node[font=\tiny,anchor=east] at (-0.03,-1.05-0.3*\rr) {\lab};
    \node[font=\scriptsize,align=center] at (1.8,-0.2) {重み行列 $\mW$（4 行 $\times$ 12 列）};
    % one-hot
    \foreach \r in {0,...,11} {
      \draw[black!45,fill=white] (3.9,-0.25*\r) rectangle ++(0.3,-0.25);
      \node[font=\tiny,anchor=west] at (4.22,-0.25*\r-0.125) {\r};}
    \draw[black!60,fill=bookblue] (3.9,-2.0) rectangle ++(0.3,-0.25);
    \node[font=\scriptsize,align=center] at (4.05,-3.4) {状態 8 の\\one-hot $\vx$};
    \node[font=\large] at (4.95,-1.5) {$=$};
    \foreach \rr in {0,...,3}
      \draw[black!60,fill=figorange] (5.35,-0.9-0.3*\rr) rectangle ++(0.4,-0.3);
    \node[font=\scriptsize,align=left,anchor=west] at (5.85,-1.5) {$Q(8,\text{上})$\\$\ \ \vdots$\\$Q(8,\text{左})$};
  \end{scope}
  \draw[->,dashed,thick,bookorange] (1.65,-2.35) to[out=-30,in=-90] node[below,font=\scriptsize,pos=0.4,text=black]{同じ 4 つの数（$\mW$ は表の転置）} (5.85,-2.1);
\end{tikzpicture}
\caption{表は、one-hot ベクトルを入れる線形層と同じもの。表 $Q$ の状態 8 の行（左の橙）と、重み行列 $\mW$ の 8 列目（右の橙。列の番号は 0 から）は、同じ 4 つの数。
one-hot ベクトル $\vx$ は 8 番目だけが 1 なので、$\mW\vx$ は $\mW$ の 8 列目を取り出す。灰色の行は、壁と終端状態の欄で使わない。
線形層を深いネットワークに取り替えれば、表に収まらない問題も扱える。}
\label{fig:rlb-table-to-net}
\end{figure}

更新の式も対応します。
目標を $y=r_t+\gamma\max_{a'}Q(s_{t+1},a')$ と置きます。
目標は「いまのところの正解の代わり」なので、教師あり学習の正解ラベルと同じく、動かさない定数として扱います。
そのうえで、損失 $\calL=\frac12\,(y-Q(s_t,a_t))^2$ を考えます。
$Q(s_t,a_t)$ で微分すると $-(y-Q(s_t,a_t))$ です。
勾配降下法（\chapref{ch:ml}）で、学習率 $\alpha$ で 1 回更新すると、次のようになります。
\[
Q(s_t,a_t)\leftarrow Q(s_t,a_t)-\alpha\cdot\bigl(-(y-Q(s_t,a_t))\bigr)=Q(s_t,a_t)+\alpha\,(y-Q(s_t,a_t))
\]
これは、Q 学習の式\eqref{eq:rlb-q-learning}と同じです。
Q 学習は、「目標との二乗誤差を、勾配降下法で 1 歩だけ減らす」ことだったのです。
PyTorch で確かめます。

\snippet{Q の表を線形層で表し、勾配降下法で更新する}{code:rlb-q-net}{ch20/q_net.py}{q-net}

1〜2 行目が、重みを 0 で始めた線形層です。
\code{net_update} は、\chapref{ch:nn}以来おなじみの学習ループの 1 回分です。
9〜10 行目で目標を計算します。\code{torch.no_grad()} の中なので、目標は定数として扱われます。
11 行目が損失で、12〜14 行目で勾配を計算して重みを更新します。
表の Q 学習と同じ経験を 1 歩ずつ、2,722 回与えて、結果を比べました。
表とネットワークの重みの差は、最大でも $1.3\times10^{-7}$ でした（\file{ch20/out/q_net.txt}）。
浮動小数点数の丸めの大きさしかありません。

線形層を深いネットワークに取り替え、Q 学習をゲームの画面に使えるようにしたのが、\term{DQN}{deep Q-network}です\cite{mnih2015human}。
49 種類のテレビゲームの多くで、人間の熟練者に並ぶ得点を出しました。
ネットワークで Q 学習をすると学習が不安定になりやすいので、過去の経験をためておいて何度も学習に使う\jterm{経験再生}などの工夫を加えています。
本書では、DQN の実装には進みません。LLM に向いているのは、次に述べる別の方法（方策ベース）だからです。

\subsection*{価値ベースと方策ベース}

この章の方法は、どれも「まず価値を学び、価値が最大の行動を選ぶ」という順で方策を作りました。
この考え方の方法を\term{価値ベース}{value-based}の方法と呼びます。
これに対して、方策そのものをパラメータを持つ関数で表し、直接良くしていく方法を、\term{方策ベース}{policy-based}の方法と呼びます（\figref{fig:rlb-map}、\tabref{tab:rlb-vb-pb}）。

\begin{figure}[tbp]
\centering
\begin{tikzpicture}
  \node[lbl,anchor=west] at (0,0.95) {\textbf{価値ベース}（この章）};
  \node[box,minimum width=21mm,minimum height=10mm] (dp) at (1.05,0) {動的計画法};
  \node[box,minimum width=27mm,minimum height=10mm] (td) at (5.15,0) {モンテカルロ法\\TD 学習};
  \node[box,minimum width=16mm,minimum height=10mm] (ql) at (8.9,0) {Q 学習};
  \node[box,minimum width=14mm,minimum height=10mm] (dqn) at (11.8,0) {DQN};
  \draw[arr] (dp) -- node[above,font=\scriptsize,align=center]{環境のモデル\\を使わず} (td);
  \draw[arr] (td) -- node[above,font=\scriptsize]{最適方策を} (ql);
  \draw[arr] (ql) -- node[above,font=\scriptsize]{表を関数に} (dqn);
  \node[lbl,anchor=west] at (0,-1.45) {\textbf{方策ベース}（次章）};
  \node[purplebox,minimum width=25mm,minimum height=10mm] (pg) at (1.25,-2.2) {方策勾配法\\{\footnotesize REINFORCE}};
  \node[purplebox,minimum width=31mm,minimum height=10mm] (ac) at (5.15,-2.2) {アクター・クリティック\\PPO};
  \node[greenbox,minimum width=40mm,minimum height=10mm] (llm) at (10.5,-2.2) {LLM の強化学習\\{\footnotesize \chapref{ch:rl-llm}、\chapref{ch:agentic-rl}}};
  \draw[arr] (pg) -- node[below,font=\scriptsize,align=center]{価値も\\補助に学ぶ} (ac);
  \draw[arr] (ac) -- node[below,font=\scriptsize,align=center]{LLM に\\当てはめる} (llm);
  \draw[darr] (td.south) -- node[right,lbl,pos=0.45]{価値関数と TD 誤差を、方策の学習の助けに使う} (ac.north);
\end{tikzpicture}
\caption{強化学習の方法の関係図。上の段は、価値を学んでから方策を作る方法で、この章で扱った。矢印に添えた語は、左の方法から右の方法へ進むときに、何を変えるか（加えるか）を示す。
下の段は、方策を直接良くする方法で、次章から扱う（下の段の名前は次章で説明する）。
下の段の右の矢印は、方法を変えるのではなく、アクター・クリティックや PPO を LLM の学習に使うことを示す。価値関数と TD 誤差は、下の段でも補助として使われる（破線）。}
\label{fig:rlb-map}
\end{figure}

\begin{table}[tbp]
\centering
\caption{価値ベースと方策ベースの比較。方策ベースの方法も、価値関数を補助に使うことが多い（\figref{fig:rlb-map}の破線）。}
\label{tab:rlb-vb-pb}
\small
\begin{tabularx}{\linewidth}{@{}l>{\raggedright\arraybackslash}X>{\raggedright\arraybackslash}X@{}}
\toprule
 & 価値ベース（この章） & 方策ベース（次章から） \\
\midrule
学習するもの & 価値 $Q(s,a)$ の表や関数 & 方策 $\pitheta(a\mid s)$ のパラメータ $\vtheta$ \\
行動の選び方 & 価値が最大の行動（$\argmax$） & 方策の確率に従ってサンプリング \\
できる方策 & 決定的（学習中は $\varepsilon$-貪欲） & 確率的なまま表せる \\
LLM に使うと & 出力の意味を「収益の見積もり」に変えて学び直す & 次のトークンの確率を、そのまま方策に使える \\
例 & Q 学習、DQN & REINFORCE、PPO（価値も補助に学ぶ） \\
\bottomrule
\end{tabularx}
\end{table}

LLM には、方策ベースが向いています（価値関数を補助に使う方法も含みます。\figref{fig:rlb-map}の破線）。
いちばんの理由は、\textbf{LLM が、はじめから方策の形をしている}ことです。
トークン列（状態）を受け取り、次のトークン（行動）の確率を出します。
しかも、事前学習と SFT を終えた、すでにかなり良い方策です。
価値ベースに切り替えると、出力の意味が「次のトークンの確率」から「そのトークンを選んだ先の収益の見積もり」に変わります。
語彙の数だけあるその値を、学び直すことになります。
すでに持っている方策を、そのまま少しずつ良くするほうが自然です。

ただし、この章で学んだことが無駄になるわけではありません。
収益、価値関数、TD 誤差は、次章で方策の学習を安定させる道具として再び登場します。
最後に、この章の実験と、実際の強化学習の規模を比べておきます（\tabref{tab:rlb-scale}）。

\begin{table}[tbp]
\centering
\caption{本書の小さな実験と、実際の強化学習の規模。AlphaGo は、囲碁で世界のトップ棋士に勝ったプログラム。DQN と AlphaGo の数値は、それぞれの論文\cite{mnih2015human,silver2016mastering}による。}
\label{tab:rlb-scale}
\footnotesize
\begin{tabularx}{\linewidth}{@{}l*{4}{>{\raggedright\arraybackslash}X}@{}}
\toprule
 & 本章の格子世界 & DQN（ゲーム） & AlphaGo（囲碁） & LLM の強化学習 \\
\midrule
状態 & マスの番号（11 個） & 画面 4 枚分の画素 & 盤面（約 $10^{170}$ 通り） & トークン列 \\
行動 & 4 個 & 最大 18 個 & 最大 361 個 & 語彙の数（数千〜数十万個） \\
価値や方策 & 36 個（9 状態 $\times$ 4 行動）の数の表 & 畳み込みネットワーク & 深いネットワークと先読み & Transformer \\
学習の量 & 2,000 エピソード（数秒） & 1 ゲームあたり 5,000 万画面 & 棋譜 3,000 万局面のあと自己対戦 & \chapref{ch:rl-llm}で扱う \\
\bottomrule
\end{tabularx}
\end{table}

\subsection*{迷路と LLM の対応}

規模だけでなく、問題の中身も比べておきます。
\tabref{tab:rlb-maze-llm}は、この章の格子世界と、\chapref{ch:rl-llm}で強化学習にかける計算アシスタントを、項目ごとに対応させたものです。
計算アシスタントでは、利用者が「12+35は？」と尋ね、LLM が「47です。」と 1 トークンずつ答えます。
この 1 回の応答が、1 エピソードに当たります。
詳しくは\secref{sec:rl-map}で、1 回の応答をトークン単位の MDP として書き下します。

\begin{table}[tbp]
\centering
\caption{格子世界と、計算アシスタントの強化学習（\chapref{ch:rl-llm}）の対応。計算アシスタントの欄は\secref{sec:rl-map}と\secref{sec:rl-exp}の設定。}
\label{tab:rlb-maze-llm}
\small
\begin{tabularx}{\linewidth}{@{}l>{\raggedright\arraybackslash}X>{\raggedright\arraybackslash}X@{}}
\toprule
 & 格子世界（この章） & 計算アシスタント（\chapref{ch:rl-llm}） \\
\midrule
目的 & 穴を避けてゴールに早く着く（期待収益 $J=V^\pi(8)$ を最大に） & 質問に正しく答える（報酬の期待値、つまり正答の確率を最大に） \\ \addlinespace
状態 & いまのマス（11 個） & 質問と、ここまでに書いたトークンの列 \\ \addlinespace
行動 & 上・右・下・左の 4 つ & 次に書くトークン 1 つ（語彙の 4,096 通り） \\ \addlinespace
遷移 & 確率的（床が滑る） & 決定的（選んだトークンを列の末尾に付ける） \\ \addlinespace
報酬 & ゴールに入れば $+1$、穴なら $-1$ & 応答を書き終えたあとに 1 回だけ（正解なら 1、不正解なら 0） \\ \addlinespace
エピソード & スタートから、ゴールか穴に着くまで & 1 つの質問への 1 回の応答（終わりの記号 \code{<|end|>} を書くまで） \\ \addlinespace
集めるデータ & 歩いた軌跡（状態・行動・報酬の列） & 1 つの質問に何本も生成させた応答と、その採点 \\ \addlinespace
学習するもの & $Q$ の表（そこから貪欲方策を作る） & 方策そのもの（LLM のパラメータ $\vtheta$。方策勾配法で更新） \\ \addlinespace
評価 & スタートの価値 $V(8)$ & 学習に使っていない問題での正答率 \\ \addlinespace
状態の再訪 & 同じマスに何度も来る & 1 つのエピソードの中では、同じ状態に二度と来ない（列が 1 歩ごとに長くなる） \\
\bottomrule
\end{tabularx}
\end{table}

表から、迷路との大きな違いが 2 つ読み取れます。
1 つ目は、遷移が決定的で、分かっていることです。
それでも動的計画法が使えないのは、状態の数が天文学的で、表を作れないからです（この節の初め）。
2 つ目は、報酬が応答の最後に 1 回だけ届くことです。
どのトークンが正解に役立ったのかを割り振る信用割当の問題（\secref{sec:rlb-what}）が、迷路よりずっと大きくなります。

この節で分かったことをまとめます。
表は one-hot ベクトルを入れる線形層と同じもので、ネットワークに取り替えれば大きな問題に進めます。
LLM ははじめから方策の形をしているので、方策を直接良くする方法を主に使います。
迷路と LLM の強化学習は、遷移が決定的なことと、状態の数が天文学的なことで、大きく違います。

\begin{summary}
\item 強化学習は、正解の代わりに報酬だけを手がかりに学びます。報酬が遅れて届くことと、自分の行動がデータを変えることが、教師あり学習との違いです。
\item 問題は、マルコフ決定過程 $(\calS,\calA,P,r,\gamma)$ として書きます。マルコフ性とは、先のことが「いまの状態と行動」だけで決まることです。
\item 確率モデルは 2 つあります。環境の遷移確率 $P(s'\mid s,a)$ は与えられたまま変わらず、学習で良くするのは方策 $\pi(a\mid s)$ です。方策が決まると、軌跡（マスの列）の確率 $p_\pi(\tau)$ が決まります。目的は、期待収益 $J(\pi)=\E[G_0]$ の最大化です。
\item 方策 $\pi(a\mid s)$ の良さは、収益 $G_t=r_t+\gamma G_{t+1}$ の期待値である価値関数 $V^\pi$、$Q^\pi$ で測ります。
\item ベルマン方程式は、いまの状態の価値を「報酬 $+$ $\gamma\times$次の状態の価値」の平均（期待方程式）や最大（最適方程式）で表します。
\item 遷移確率が分かっていれば、方策反復や価値反復で最適方策が求まります。分からなくても、モンテカルロ法や TD 学習（TD 誤差 $\delta_t=r_t+\gamma V(s_{t+1})-V(s_t)$ で 1 歩ごとに直す）で、経験から価値を学べます。
\item Q 学習は、$\varepsilon$-貪欲法で探索しながら行動価値の表を学び、環境のモデルなしで最適方策を見つけます。
\item 表は大きな問題には使えません。LLM ははじめから方策の形をしているので、次章の方策ベースの方法を主に使います。
\end{summary}

\section*{次の章へ}

この章では、価値の表を学び、そこから方策を作りました。
しかし、LLM は方策そのものです。
それなら、価値を経由せずに、方策のパラメータを直接動かして、良い結果を増やせばよいはずです。
ところが、報酬は「サンプリングした行動の結果」として届くので、そのままでは微分できません。
計算アシスタントが「57です。」と答えて返ってくるのは、「不正解」という判定 1 つです。
この判定は、パラメータの式になっていないので、誤差逆伝播で勾配を求められません。
次の\chapref{ch:policy-gradient}では、この問題を解決する方策勾配法を導き、REINFORCE から PPO までを実装します。

\section*{さらに学ぶために}

\begin{itemize}
\item Sutton・Barto『Reinforcement Learning: An Introduction』\cite{sutton2018rl}: 強化学習の標準的な教科書。第 2〜6 章が本章の内容に当たる。全文がウェブで公開されている。
\item Watkins・Dayan「Q-learning」\cite{watkins1992q}: Q 学習が $Q^\ast$ に収束することを証明した論文。
\item Mnih ほか「DQN」\cite{mnih2015human}: Q 学習と深いネットワークを組み合わせ、画面だけからゲームを学習した論文。
\item Achiam「Spinning Up in Deep RL」\cite{achiam2018spinningup}: 深層強化学習の入門サイト。次章の予習にもなる。
\end{itemize}

\begin{exercises}
\exfirst[5 は、\chapref{ch:rl-llm}への準備になります。]{1 と 5}
\item （手計算）ランダム方策の価値 $V^\pi(2)=0.2616$、$V^\pi(6)=-0.3371$ を使って、式\eqref{eq:rlb-q-from-v}から $Q^\pi(6,\text{上})$ を計算してください。
状態 6 で「上」を選んだときの 3 通りの結果は、\secref{sec:rlb-mdp}の実行結果にあります（答えは $0.0580$）。
\item （手計算）状態が 1 つだけで、どの行動を選んでも報酬 1 をもらって同じ状態に戻る、終わりのない MDP を考えます。
ベルマン期待方程式から、$V=1/(1-\gamma)$ となることを示してください。
\item （コード）割引率 \code{GAMMA}（\file{ch20/policy_eval.py}）を 0.5 に変えて、価値反復を実行してください。
最適方策はどの状態で変わりますか。
次に、割引率を 0.9 に戻し、\file{ch20/gridworld.py} の \code{reward} を書き換えて、ゴールと穴以外のマスに入ったときの報酬を $-0.04$（1 歩ごとの小さな罰）にしてください。
最適方策はどう変わりますか。それぞれ理由も考えてください（ヒント: 穴に近い状態 9 と 11 の矢印に注目）。
\item （コード）Q 学習の目標 $r_t+\gamma\max_{a'}Q(s_{t+1},a')$ を、$r_t+\gamma Q(s_{t+1},a_{t+1})$ に変えてください（SARSA と呼ばれる方法です）。
$a_{t+1}$ は、次の状態で $\varepsilon$-貪欲法が実際に選んだ行動です。学習した $Q$ と貪欲方策を、Q 学習の結果と比べてください。
\item （考察）計算アシスタントが「12+35は？」に 1 トークンずつ答える過程を MDP と見ると、状態、行動、遷移確率、報酬、エピソードの終わりは、それぞれ何に当たるでしょうか。
\end{exercises}
